\chapter{Conormal distributions}\label{ch:conormal}
\providecommand{\osc}[1]{\mathcal{I}[#1]}
\providecommand{\IH}[1]{IH^{#1}}
\providecommand{\Vt}{\mathcal{V}}
\providecommand{\Ffib}{\mathcal{F}_{\mathrm{fib}}}

After smooth functions, the simplest distributions are those whose
singularities are as regular as possible: they sit on a smooth submanifold
$Y$, point only in the directions normal to $Y$, and look the same at all
small scales up to a power. In \cref{ch:microlocal} we met examples: the
Heaviside function $H(x_1)$, the delta function and the kernel of a
pseudodifferential operator, whose wavefront sets lie in the conormal bundle
of a submanifold (\cref{microlocal:fig:wf-examples,microlocal:fig:kernel}).
This chapter builds the class of all such distributions.

The Dirac delta function, the Schwartz kernels of pseudodifferential operators
and the travelling-wave solutions of hyperbolic equations are all examples of
distributions that are smooth outside a submanifold $Y$ and have a very
particular, ``symbolic'', singularity at $Y$.  In this chapter we introduce the
corresponding spaces $\Ic{m}(X,Y)$ of \emph{conormal distributions}.  We start
with the simplest case, $Y$ a point in $\R^n$, where conormal distributions are
simply inverse Fourier transforms of symbols; pass to a linear subspace
$Y=\{0\}\times\R^d\subset\R^n$ (equivalently the zero section of a trivial
vector bundle), where the Fourier transform is taken only in the normal
variables; prove the invariance of these spaces under diffeomorphisms preserving
$Y$; and then define $\Ic{m}(X,Y;\Omegah)$ for a closed embedded submanifold of
a manifold, either through coordinates or through a normal fibration.  The
principal symbol is an invariantly defined half-density on the conormal bundle
$N^*Y$, and gives a short exact sequence.  Finally we show that the conormal
distributions are characterized, up to an arbitrarily small loss in the order,
by \emph{iterated regularity}: stability of a fixed Sobolev regularity under
repeated application of vector fields tangent to $Y$.

We use the symbol spaces of \cref{ch:euclidean}, the radial compactification
of \cref{ch:distributions} and, for wavefront sets, \cref{ch:microlocal}.
Pseudodifferential operators on manifolds are discussed only in
\cref{ch:manifolds}; their action on conormal distributions is derived here in
local coordinates and globalized in \cref{ch:hj}.

\section{Symbols}\label{conormal:sec:symbols}

Conormal distributions will be Fourier transforms of symbols in the normal
variables, so we first collect what we need about symbols, and in particular
about classical symbols, which are smooth functions on the radial
compactification.

\subsection{Symbols with bounds}
Recall from \cref{ch:euclidean} that for $m\in\R$, $S^m(\R^p;\R^k)$ is the
space of $a\in\CI(\R^p\times\R^k)$ such that
\begin{equation}\label{conormal:eq:symbol-est}
  |D_z^\alpha D_\xi^\beta a(z,\xi)|\le C_{\alpha\beta}\ang{\xi}^{m-|\beta|}
  \quad\forall\ \alpha,\beta,
\end{equation}
with the Fr\'echet topology given by the best constants; $S^m(\R^k)$ denotes the
case without parameters, $S^{-\infty}=\bigcap_m S^m$, and every sequence
$a_j\in S^{m-j}$ has an asymptotic sum $a\sim\sum_j a_j$, unique modulo
$S^{-\infty}$.  We shall also write $S^m_{\mathrm c}(\R^p;\R^k)$ for the
subspace of symbols vanishing for $z$ outside a compact subset of $\R^p$.
Directly from \eqref{conormal:eq:symbol-est}, $\partial_{\xi_j}:S^m\to S^{m-1}$, $\xi_j:S^m\to S^{m+1}$, and
$S^m\cdot S^{m'}\subset S^{m+m'}$.  Note that $S^{-\infty}(\R^k)=\cS(\R^k)$.

The symbol spaces are invariant under linear changes of the variable: if
$A(z)\in\GL(k,\R)$ depends smoothly on $z$, with $A$, $A^{-1}$ and all their
derivatives bounded, then $a(z,A(z)\xi)\in S^m(\R^p;\R^k)$ whenever
$a\in S^m(\R^p;\R^k)$, since $\ang{A(z)\xi}$ and $\ang{\xi}$ are uniformly
comparable and each $z$-derivative falling on $A(z)\xi$ produces a factor of
order one together with a $\xi$-derivative of $a$.  In particular the symbol
spaces $S^m(W)$ are well defined for any finite-dimensional real vector space
$W$.

The following density statement allows many identities to be proved first for
Schwartz functions.

\begin{lemma}\label{conormal:lem:density}
  If $a\in S^m(\R^p;\R^k)$, $\phi\in\CIc(\R^k)$ equals $1$ near $0$ and
  $a_\epsilon(z,\xi)=\phi(\epsilon\xi)a(z,\xi)$, then $a_\epsilon$ is bounded in
  $S^m$ for $0<\epsilon\le1$ and $a_\epsilon\to a$ in $S^{m'}$ for every $m'>m$
  as $\epsilon\to0$.
\end{lemma}
\begin{proof}
  The derivative $D_\xi^\gamma[\phi(\epsilon\xi)]=\epsilon^{|\gamma|}
  (D^\gamma\phi)(\epsilon\xi)$ is supported where $|\xi|\le C/\epsilon$, so it is
  bounded by $C_\gamma\ang{\xi}^{-|\gamma|}$ uniformly in $\epsilon$; hence
  $a_\epsilon$ is bounded in $S^m$ by Leibniz' formula.  Moreover
  $a-a_\epsilon=(1-\phi(\epsilon\xi))a$ and its derivatives are supported in
  $|\xi|\ge c/\epsilon$, where $\ang{\xi}^{m-|\beta|}\le
  C\epsilon^{m'-m}\ang{\xi}^{m'-|\beta|}$; so $a-a_\epsilon\to0$ in $S^{m'}$.
\end{proof}

\subsection{Classical symbols}
Let $W$ be a real vector space of dimension $k$ and $\RC{W}$ its radial
compactification (\cref{ch:distributions}), a compact manifold with boundary
whose interior is $W$; the action of $\GL(W)$ and of translations extends
smoothly to $\RC W$.  We call the boundary $\Sphere W=\pa\RC W$ the
\emph{sphere at infinity}; it is canonically $(W\setminus 0)/\R^+$.  Let
$\rho\in\CI(\RC W)$ be a boundary defining function, for instance
$\rho=\ang{\xi}^{-1}$ in linear coordinates.  Any two are related by
$\rho'=f\rho$ with $0<f\in\CI(\RC W)$, so
\begin{equation*}
  \Symph{m}(W)=\rho^{-m}\CI(\RC W)=\{a\in\CI(W);\ \rho^m a\in\CI(\RC W)\}
\end{equation*}
is independent of the choice; this is the space of \emph{classical} (or
polyhomogeneous) symbols of order $m$.  (One may regard it as the space of
sections of a trivial line bundle $N_{-m}$ over $\RC W$ of which $\rho^{-m}$ is
a global trivializing section.)

\begin{discussion}
  Why compactify? A classical symbol has an expansion in homogeneous terms of
  decreasing degree at infinity. After the substitution $r=\abs\xi^{-1}$ this
  is simply a Taylor expansion in $r$ at $r=0$, the sphere at infinity
  (\cref{conormal:fig:compactify}). So statements about expansions, such as
  asymptotic summation and the principal part, become statements about smooth
  functions on a manifold with boundary and their Taylor series at the
  boundary. The one-point compactification would collapse the sphere at
  infinity to a point and forget the direction $\xi/\abs\xi$, which is
  exactly where the principal part lives.
\end{discussion}

\begin{lemma}\label{conormal:lem:classical-basic}
  For all $m,m'\in\R$:
  \begin{enumerate}[label=(\roman*)]
  \item $\Symph{m}(W)\cdot\Symph{m'}(W)=\Symph{m+m'}(W)$ and $\Symph{m}(W)$ is
    invariant under $\GL(W)$ and translations;
  \item $\Symph{m}(W)\subset S^m(W)$; in linear coordinates
    $\xi^\alpha\pa_\xi^\beta:\Symph{m}\to\Symph{m+|\alpha|-|\beta|}$;
  \item $\bigcap_N\Symph{m-N}(W)=\cS(W)$;
  \item $a\in\Symph{m}(W)$ if and only if there are functions
    $a_{m-j}\in\CI(W\setminus0)$, homogeneous of degree $m-j$, such that for
    one (hence any) $\chi\in\CIc(W)$ equal to one near $0$,
    \begin{equation}\label{conormal:eq:ph-expansion}
      a-\sum_{j<N}(1-\chi)a_{m-j}\in S^{m-N}(W)\quad\forall\ N.
    \end{equation}
  \end{enumerate}
\end{lemma}
\begin{proof}
  (i) is immediate from the definition and the invariance of $\RC W$.  For
  (ii) and (iv) use linear coordinates and, near the boundary, the coordinates
  $r=|\xi|^{-1}\in[0,1)$, $\omega=\xi/|\xi|\in\Sphere^{k-1}$ on $\RC W$; then
  $\rho=r\,g$ with $g>0$ smooth up to $r=0$, so $a\in\Symph{m}$ means
  $a=r^{-m}f(r,\omega)$ with $f$ smooth up to $r=0$.  In these coordinates
  $\xi_j\pa_{\xi_l}$ is a combination, with coefficients smooth up to $r=0$, of
  $r\pa_r$ and vector fields on $\Sphere^{k-1}$, while $\pa_{\xi_l}$ is $r$ times
  such a combination,
  as one sees from homogeneity ($\xi_j\pa_{\xi_l}$ is homogeneous of degree $0$
  and $\pa_{\xi_l}$ of degree $-1$).  Since $r\pa_r(r^{-m}f)=r^{-m}(r\pa_r-m)f$,
  repeated application of such vector fields to $r^{-m}f$ yields $r^{-m}$ times
  bounded functions, so the estimates \eqref{conormal:eq:symbol-est} hold in
  $|\xi|\ge1$ (where $r\sim\ang{\xi}^{-1}$), and trivially in $|\xi|\le1$; the
  mapping property in (ii) follows as well.  Taylor's formula in $r$ at $r=0$,
  $f=\sum_{j<N}r^jf_j(\omega)+r^Nf_N'(r,\omega)$, gives
  \eqref{conormal:eq:ph-expansion} with
  $a_{m-j}(\xi)=|\xi|^{m-j}f_j(\xi/|\xi|)$, the remainder $r^{N-m}f'_N$ being in
  $S^{m-N}$ by the argument just given.  Conversely, if
  \eqref{conormal:eq:ph-expansion} holds, then near $r=0$ the function
  $f=r^{m}a$ satisfies $f-\sum_{j<N}r^jf_j=O(r^N)$ together with all
  derivatives in $r\pa_r$ and $\pa_\omega$, and $\pa_r=r^{-1}(r\pa_r)$; so $f$ is
  $\mathcal C^{N-1}$ up to $r=0$ for every $N$, i.e.\ smooth.  Finally (iii): the
  intersection is contained in $S^{-\infty}(W)=\cS(W)$ by (ii), and conversely a
  Schwartz function is $r^\infty$ times a smooth function of $(r,\omega)$, by
  the same change of variables.
\end{proof}

The leading term in \eqref{conormal:eq:ph-expansion} is the \emph{principal
part} $\sigma_m(a)=a_m$; it is a smooth function on $W\setminus0$ homogeneous
of degree $m$, equivalently (multiplying by $\rho^{m}$ and restricting) a
smooth section of $N_{-m}$ over $\Sphere W$.  Denote by
$\CI_{(m)}(W\setminus0)$ the space of smooth functions homogeneous of degree $m$.

\begin{proposition}\label{conormal:prop:ph-exact}
  For every $m$ the sequence
  \begin{equation*}
    0\longrightarrow\Symph{m-1}(W)\longrightarrow\Symph{m}(W)
    \xrightarrow{\ \sigma_m\ }\CI_{(m)}(W\setminus0)\longrightarrow0
  \end{equation*}
  is exact.  Equivalently $\rho\CI(\RC W)\to\CI(\RC W)\to\CI(\Sphere W)$ is
  exact: restriction to the boundary is surjective and a smooth function
  vanishing at the boundary is $\rho$ times a smooth function.
\end{proposition}
\begin{proof}
  Multiplication by $\rho^m$ reduces the statement to the case $m=0$, where it
  is Taylor's theorem at the boundary (in the coordinates $(r,\omega)$ above) and
  the existence of extensions: $f(\omega)\mapsto(1-\chi)f(\xi/|\xi|)$.
\end{proof}

\begin{proposition}[Asymptotic completeness]\label{conormal:prop:asymp}
  If $a_j\in\Symph{m-j}(W)$, $j\in\N_0$, there exists $a\in\Symph{m}(W)$ such
  that $a-\sum_{j<N}a_j\in\Symph{m-N}(W)$ for all $N$; $a$ is unique modulo
  $\cS(W)$.  The same holds, with the same proof, for symbols depending
  smoothly on parameters in a compact set.
\end{proposition}
\begin{proof}
  Uniqueness follows from \cref{conormal:lem:classical-basic}(iii).  Multiplying
  by $\rho^m$, it suffices to show: if $X$ is a compact manifold with boundary,
  $\rho$ a boundary defining function and $b_j\in\rho^j\CI(X)$, there is
  $b\in\CI(X)$ with $b-\sum_{j<N}b_j\in\rho^N\CI(X)$ for all $N$.  Using a
  collar $[0,1)_r\times\pa X$ and a partition of unity this reduces to the
  following form of Borel's lemma: given $u_j\in\CI(\pa X)$ there is
  $u\in\CIc([0,1)\times\pa X)$ with $\pa_r^ju(0,\cdot)=u_j$ for all $j$.  (Indeed,
  write $b_j=\sum_{l}r^{j+l}c_{jl}(\omega)/(j+l)!+O(r^\infty)$ near the boundary and
  take the $u$ with $u_p=\sum_{j\le p}c_{j,p-j}$; then $b-\sum_{j<N}b_j$
  vanishes to order $N$ at $r=0$.)  Choose $\chi\in\CIc(\R)$ with $\chi=1$ on
  $|r|<\frac12$, $\chi=0$ on $|r|>1$, and set
  \begin{equation*}
    u(r,\omega)=\sum_j \frac{r^j}{j!}u_j(\omega)\chi(r/\epsilon_j),\qquad
    0<\epsilon_j\le1.
  \end{equation*}
  Each term vanishes for $r>\epsilon_j$, and for $l\le j$ and any differential
  operator $Q$ of order $q$ on $\pa X$,
  $|\pa_r^l Q(r^ju_j\chi(r/\epsilon_j)/j!)|\le C_{j,l,Q}\,\epsilon_j^{j-l}$,
  with $C_{j,l,Q}$ independent of $\epsilon_j$, since
  $\pa_r^p[\chi(r/\epsilon_j)]=\epsilon_j^{-p}\chi^{(p)}(r/\epsilon_j)$ and
  $r\le\epsilon_j$ on the support.  Fix an exhausting sequence
  $Q_1,Q_2,\dots$ of differential operators on $\pa X$.  Choosing
  $\epsilon_j\le 2^{-j}/(1+\max_{l+q\le j-1}C_{j,l,Q_q})$ imposes only finitely
  many conditions on each $\epsilon_j$, and makes every differentiated series
  converge uniformly.  Hence $u$ is smooth and, since only the term $j$
  contributes to $\pa_r^j u(0,\cdot)$, $\pa_r^ju(0,\cdot)=u_j$.  Parameters are
  handled by including the parameter derivatives among the $Q_q$.
\end{proof}

The symbol estimates \eqref{conormal:eq:symbol-est} do not imply that $a$ is
classical: $\Symph{m}(W)$ is a closed subspace of $S^m(W)$ in the topology of
$\rho^{-m}\CI(\RC W)$, and $\cS(W)$ is not dense in it in that topology, though
by \cref{conormal:lem:density} it is dense in the weaker topology of
$S^{m'}$, $m'>m$.  The relation between the two classes is explored in
\cref{conormal:ex:Vkl,conormal:ex:Ab}.  For symbols depending on parameters
$z\in\R^p$ we write $\Symph{m}(\R^p;\R^k)$ for those $a\in S^m(\R^p;\R^k)$
having an expansion \eqref{conormal:eq:ph-expansion} with
$a_{m-j}\in\CI(\R^p\times(\R^k\setminus0))$ homogeneous in $\xi$, uniformly
bounded with all derivatives on $\R^p\times\Sphere^{k-1}$, the remainders being
in $S^{m-N}(\R^p;\R^k)$.  All statements in this chapter hold for the classical
spaces as well as for the spaces of symbols with bounds; we indicate the
classical versions by a subscript ``$\mathrm{ph}$''.

\section{Conormal distributions at a point}\label{conormal:sec:point}

The model is the delta function: $\hat\delta_0=1$ is the simplest symbol. We
allow any symbol on the Fourier side.

\begin{definition}\label{conormal:def:point}
  For $m\in\R$ set
  \begin{equation*}
    \Ic{m}_{\cS}(\R^n,\{0\})=\{u\in\cSp(\R^n);\ \hat u\in S^{m-\frac n4}(\R^n)\},
    \qquad
    \Ic{m}_{\cS,\mathrm{ph}}(\R^n,\{0\})=\cF^{-1}\big(\Symph{m-\frac n4}(\R^n)\big).
  \end{equation*}
\end{definition}

The shift $n/4$ in the order is a normalization, which is explained by the
behaviour of the order under the operations considered later (it makes the
order of the Schwartz kernel of a pseudodifferential operator of order $m$, as
a conormal distribution at the diagonal, equal to $m$).  For example
$\hat\delta_0=1$, so
\begin{equation*}
  \delta_0\in\Ic{\frac n4}_{\cS,\mathrm{ph}}(\R^n,\{0\}).
\end{equation*}
Since $S^{m-1}\subset S^m$, $\Ic{m-1}_\cS\subset\Ic{m}_\cS$, and
$\bigcap_m\Ic{m}_\cS(\R^n,\{0\})=\cF^{-1}\cS(\R^n)=\cS(\R^n)$.

\begin{question}
  Why ask $\hat u$ to be a symbol, and not just polynomially bounded?
\end{question}
\begin{answer}
  Every $u\in\CmIc(\R^n)$ has a polynomially bounded Fourier transform, so that
  alone says nothing about the singularity. The extra information in the
  symbol estimates is that each $\zeta$-derivative gains a power of
  $\ang\zeta^{-1}$. Since $D_\zeta$ corresponds to multiplication by $-z$, this
  says that $z^\alpha u$ becomes smoother as $\abs\alpha$ grows: the singularity
  is concentrated at $z=0$ in a uniform, scale-invariant way. This is the
  content of \cref{conormal:lem:point-basic}(i)--(iii) below.
\end{answer}

\begin{lemma}\label{conormal:lem:point-basic}
  Let $u\in\Ic{m}_\cS(\R^n,\{0\})$ (resp.\ the classical space).  Then
  \begin{enumerate}[label=(\roman*)]
  \item $z^\alpha D_z^\beta u\in\Ic{m+|\beta|-|\alpha|}_\cS(\R^n,\{0\})$ (resp.\
    classical);
  \item if $m<-\frac{3n}4-j$, $j\in\N_0$, then $u\in\mathcal C^j(\R^n)$ and
    $D^\beta u\in\Czinf(\R^n)$ for $|\beta|\le j$;
  \item if $\chi\in\CIc(\R^n)$ equals one near $0$, then $(1-\chi)u\in\cS(\R^n)$;
    in particular $\singsupp u\subset\{0\}$;
  \item the integral $\int u\,dz$, defined as $\int\chi u+\int(1-\chi)u$
    (independent of $\chi$), equals $\hat u(0)$;
  \item if $v\in\Ic{m'}_\cS(\R^n,\{0\})$ then $u*v\in
    \Ic{m+m'-\frac n4}_\cS(\R^n,\{0\})$.
  \end{enumerate}
\end{lemma}
\begin{proof}
  (i) follows from $\cF(z^\alpha D^\beta u)=(-D_\zeta)^\alpha(\zeta^\beta\hat
  u)$ and \cref{conormal:lem:classical-basic}(ii).  (ii): if $m<-3n/4-j$ then
  $\zeta^\beta\hat u\in S^{m-\frac n4+|\beta|}\subset L^1$ for $|\beta|\le j$,
  so $D^\beta u$ is the inverse Fourier transform of an integrable function,
  hence continuous and vanishing at infinity by the Riemann--Lebesgue lemma.
  (iii): by (i) and (ii), $|z|^{2N}D^\beta u=\sum_{|\gamma|=N}c_\gamma
  z^{2\gamma}D^\beta u$ is bounded if $2N>m+|\beta|+3n/4$; since $1-\chi$ vanishes
  near $0$, $|D^\beta((1-\chi)u)|\le C_{N,\beta}\ang{z}^{-2N}$ for all $N$ and
  $\beta$.  (iv): the definition is independent of $\chi$ since the two
  definitions agree on $\CIc(\R^n)$; and $\int\chi u=\widehat{\chi u}(0)$,
  $\int(1-\chi)u=\widehat{(1-\chi)u}(0)$ add up to $\hat u(0)$.  (v): $\widehat{u*v}=
  \hat u\hat v\in S^{m-\frac n4+m'-\frac n4}=S^{(m+m'-\frac n4)-\frac n4}$.
\end{proof}

\subsection{Densities and invariance}
To see that these spaces are naturally associated to the origin of an abstract
vector space we need densities.  Let $W$ be a real vector space of dimension
$n$.  For $t\in\R$, the space $\Omega^tW$ of \emph{$t$-densities} consists of the
maps $\nu:\Lambda^nW\setminus0\to\R$ with $\nu(cv)=|c|^t\nu(v)$ for
$c\in\R\setminus0$; it is a line, and $\Omega W=\Omega^1W$.  If $G\in\GL(W)$
then $G^*\nu=|\det G|^t\nu$; for $W=\R^n$ we write $|dz|^t$ for the $t$-density
taking the value $1$ on $e_1\wedge\dots\wedge e_n$.  Densities are what one
integrates invariantly: $\int G^*\mu=\int\mu$ for $\mu\in\cS(W;\Omega W)$.  The
pairing $\Omega W\otimes\Omega W'\to\R$, $|dz|\otimes|d\zeta|\mapsto1$ for dual
bases, is independent of the basis, so $\Omega W'=(\Omega W)'$.

For $u\in\cS(\R^n)$ and $G\in\GL(n,\R)$ the change of variables $y=Gz$ gives
\begin{equation*}
  \widehat{G^*u}(\zeta)=\int e^{-iz\cdot\zeta}u(Gz)\,dz
  =|\det G|^{-1}\hat u\big((G^{-1})^t\zeta\big).
\end{equation*}
Thus if we define the Fourier transform to take values in densities,
$\cF u=\hat u(\zeta)|d\zeta|$, then $\cF(G^*u)=((G^{-1})^t)^*\cF u$: the
Fourier transform is a linear-invariant isomorphism
\begin{equation}\label{conormal:eq:invariant-FT}
  \cF:\cS(W)\longrightarrow\cS(W';\Omega W'),\qquad
  \cF^{-1}\mu(w)=(2\pi)^{-n}\int_{W'}e^{i\pair w\zeta}\mu(\zeta),
\end{equation}
extending to $\cSp(W)\to\cSp(W';\Omega W')$.  Consequently
\begin{equation*}
  \Ic{m}(W,\{0\})=\cF^{-1}\big(S^{m-\frac n4}(W';\Omega W')\big),
\end{equation*}
where $S^r(W';\Omega W')=S^r(W')\cdot|d\zeta|$, is well defined independently of
any choice of basis, and similarly for the classical space; the integral of
\cref{conormal:lem:point-basic}(iv) becomes the invariant map
$\Ic{m}(W,\{0\};\Omega W)\ni u\mapsto(\cF u)(0)$ if $u$ is density-valued.  The
\emph{principal symbol} of $u$ is the class of $\cF u$ in
$S^{m-\frac n4}(W';\Omega W')/S^{m-\frac n4-1}(W';\Omega W')$ and, in the
classical case, its principal part, a density on $W'\setminus0$ homogeneous of
degree $m-\frac n4+n=m+\frac{3n}{4}$ (the density factor $|d\zeta|$ being
homogeneous of degree $n$).  By \cref{conormal:prop:ph-exact} there are short
exact sequences
\begin{equation*}
  \Ic{m-1}(W,\{0\})\hookrightarrow\Ic{m}(W,\{0\})\xrightarrow{\ \sigma_m\ }
  S^{m-\frac n4}(W';\Omega W')/S^{m-\frac n4-1}(W';\Omega W'),
\end{equation*}
and, in the classical case, onto the homogeneous densities of degree
$m+\frac{3n}4$ on $W'\setminus0$.

Invariance under non-linear changes of coordinates fixing $0$ is a special case
of \cref{conormal:thm:invariance} below.

\subsection{Homogeneity and conormality}
For classical conormal distributions the singularity at the origin can be
described quite explicitly.  We use two facts about homogeneous distributions
\cite[Theorems~3.2.3, 7.1.16, 7.1.18]{Hormander-I}: if $f\in\CI(\R^n\setminus0)$
is homogeneous of degree $\mu$ and $\mu\notin\{-n,-n-1,\dots\}$, then $f$ has a
unique extension to a distribution $\tilde f\in\cSp(\R^n)$ homogeneous of degree
$\mu$; and the Fourier transform of a homogeneous distribution of degree $\mu$
which is smooth outside $0$ is homogeneous of degree $-n-\mu$ and smooth outside
$0$.

\begin{proposition}\label{conormal:prop:homogeneous}
  Let $u\in\Ic{m}_{\cS,\mathrm{ph}}(\R^n,\{0\})$, so $\hat u\in\Symph{\mu}(\R^n)$
  with $\mu=m-\frac n4$, and suppose $\mu\notin\Z$.  Then there are functions
  $U_j\in\CI(\R^n\setminus0)$, homogeneous of degree $-n-\mu+j$ and determined by
  the homogeneous terms $\hat u_{\mu-j}$ in the expansion of $\hat u$, such that
  \begin{equation*}
    u-\sum_{j<N}U_j\in\mathcal C^r(\R^n)\quad\text{for } 0\le r<N-\mu-n.
  \end{equation*}
  Conversely the $U_j$ determine the terms $\hat u_{\mu-j}$, so $u$ is determined
  modulo $\cS(\R^n)$ by its expansion.
\end{proposition}
\begin{proof}
  Let $\chi\in\CIc(\R^n)$ equal one near $0$.  Since $\mu-j$ is not an integer,
  $\hat u_{\mu-j}$ has a homogeneous extension $\tilde a_j$, and
  $U_j:=\cF^{-1}\tilde a_j$ is homogeneous of degree $-n-\mu+j$ and smooth
  outside $0$.  The difference $(1-\chi)\hat u_{\mu-j}-\tilde a_j=-\chi\tilde a_j$
  has compact support, so its inverse Fourier transform is smooth.  The remainder
  $\hat u-\sum_{j<N}(1-\chi)\hat u_{\mu-j}\in S^{\mu-N}$ has inverse Fourier
  transform in $\mathcal C^r$ if $\mu-N+r<-n$.  The converse follows by
  applying the Fourier transform to the homogeneous terms, which are
  distinguished by their degrees.
\end{proof}

When $\mu$ is an integer the expansion of $\hat u$ is unchanged, but the
expansion of $u$ may contain terms $p(z)\log|z|$ with $p$ a homogeneous
polynomial (for $\mu\le-n$), coming from the homogeneous terms of degree
$-n-l$, $l\in\N_0$, which have no homogeneous extension; and for
$\mu\in\N_0$ the polynomial part of $\hat u$ corresponds to derivatives of the
delta function at $0$, invisible in $z\neq0$.

\subsection{Blow-up of the origin}
A construction dual to radial compactification is the blow-up of the origin,
which resolves the singularity of such homogeneous functions.  Polar coordinates
give a smooth surjective map
\begin{equation*}
  \beta:[0,\infty)\times\Sphere^{n-1}\ni(r,\omega)\longmapsto r\omega\in\R^n,
\end{equation*}
a diffeomorphism of $(0,\infty)\times\Sphere^{n-1}$ onto $\R^n\setminus0$ which
collapses $\{0\}\times\Sphere^{n-1}$ to $0$.  Each $U_j$ in
\cref{conormal:prop:homogeneous} is of the form $r^{-n-\mu+j}$ times a smooth
function on the sphere.  Radial compactification adds a sphere at infinity;
blowing up replaces the origin by a sphere (\cref{conormal:fig:compactify}).
In both cases a homogeneous function becomes a power of the boundary defining
function times a smooth function.

\begin{figure}[htbp]
\centering
\begin{tikzpicture}[>=Latex,scale=0.9]
  \begin{scope}
    \fill[gray!12] (0,0) circle (2);
    \draw[very thick] (0,0) circle (2);
    \foreach \a in {15,75,...,375} \draw[->] (0,0)--(\a:1.95);
    \fill[gray!35] (40:1.55) arc (40:70:1.55)--(70:2) arc (70:40:2)--cycle;
    \node[font=\footnotesize] at (55:2.45) {$r<\epsilon$};
    \node[font=\footnotesize,anchor=west] at (-30:2.1) {$\Sphere W=\{r=0\}$};
    \node[font=\footnotesize] at (0,-2.6) {(a) radial compactification $\RC W$};
  \end{scope}
  \begin{scope}[xshift=7.2cm]
    \fill[gray!12] (0,0) circle (2);
    \fill[white] (0,0) circle (0.7);
    \draw[very thick] (0,0) circle (0.7);
    \draw[dashed] (0,0) circle (2);
    \foreach \a in {15,75,...,375} \draw[->] (\a:0.7)--(\a:1.9);
    \node[font=\footnotesize,align=center] at (0,0) {$\{r=0\}$\\$=\Sphere^{n-1}$};
    \node[font=\footnotesize] at (0,-2.6) {(b) blow-up $[W,\{0\}]$};
    \draw[->] (2.3,0)--(3.3,0) node[midway,above,font=\footnotesize]{$\beta$};
    \fill (3.8,0) circle (1.6pt) node[below,font=\footnotesize]{$0\in W$};
  \end{scope}
\end{tikzpicture}
\caption[Radial compactification and blow-up]{(a) The radial
compactification $\RC W$ adds the sphere at infinity $\Sphere W$, with
boundary defining function $r=\abs\xi^{-1}$ near it. Rays become arcs ending at
the boundary, and a neighbourhood of a boundary point (dark) is a conic
region $\{\abs\xi>1/\epsilon\}$ around a direction. Classical symbols of order
$m$ are $r^{-m}$ times smooth functions on $\RC W$. (b) The blow-up of the
origin replaces $0$ by the sphere $\{r=0\}$ of directions, with
$r=\abs z$; the blow-down map $\beta$ collapses it to $0$. The terms $U_j$ of
\cref{conormal:prop:homogeneous} are powers of $r$ times smooth functions on
$[W,\{0\}]$.}
\label{conormal:fig:compactify}
\end{figure}

\begin{lemma}\label{conormal:lem:blowup}
  The action of $\GL(n,\R)$ on $\R^n$ lifts under $\beta$ to a smooth action on
  $[0,\infty)\times\Sphere^{n-1}$.
\end{lemma}
\begin{proof}
  For $G\in\GL(n,\R)$ the lift is
  $\tilde G(r,\omega)=(r|G\omega|,\,G\omega/|G\omega|)$, which is smooth up to
  $r=0$, satisfies $\beta\circ\tilde G=G\circ\beta$, and $\widetilde{GH}=\tilde
  G\tilde H$.  (Infinitesimally: the linear vector fields $z_i\pa_{z_j}$ are
  homogeneous of degree $0$, so on $(0,\infty)\times\Sphere^{n-1}$ they are of the
  form $a(\omega)r\pa_r+V(\omega)$ with $V$ a vector field on the sphere, which
  extends smoothly to $r=0$, tangent to the boundary.)
\end{proof}
Thus for a vector space $W$ the blown-up space $[W,\{0\}]$, a manifold with
boundary $\Sphere W$ with a blow-down map onto $W$, is well defined; see
\cref{conormal:ex:blowup}.

\section{The local model}\label{conormal:sec:local}

For a submanifold $Y$ of positive dimension we take the Fourier transform only
in the variables normal to $Y$, and let the symbol depend smoothly on the
point of $Y$. Thus a conormal distribution is, roughly, a smooth family,
parametrized by $Y$, of conormal distributions at a point.

\subsection{Notation}
Throughout the rest of the chapter $n$ is the dimension of the ambient space,
$k$ the codimension of the submanifold and $d=n-k$ its dimension.  In the local
model
\begin{equation*}
  \R^n=\R^k_x\times\R^d_y,\qquad Y=\{x=0\},\qquad
  z=(x,y),\ \zeta=(\xi,\eta),
\end{equation*}
and the order of a conormal distribution is related to the order of its
symbol by
\begin{equation}\label{conormal:eq:order}
  M=m+\frac n4-\frac k2 .
\end{equation}
For a point ($k=n$) this is $M=m-\frac n4$ as in \cref{conormal:def:point};
for the diagonal in $\R^n\times\R^n$ ($n\to2n$, $k=n$) it is $M=m$.

For an amplitude $c\in S^M(\R^n;\R^k)$, i.e.\ $c=c(x,y,\xi)$ satisfying the
estimates \eqref{conormal:eq:symbol-est} uniformly in $(x,y)$, we define the
distribution
\begin{equation}\label{conormal:eq:osc}
  \osc{c}(x,y)=\int_{\R^k}e^{ix\cdot\xi}c(x,y,\xi)\,\dbar\xi,
  \qquad
  \pair{\osc{c}}{\phi}=\int\Big(\int e^{ix\cdot\xi}c(x,y,\xi)\phi(x,y)\,dx\,dy\Big)
  \dbar\xi,
\end{equation}
for $\phi\in\CIc(\R^n)$; here $\dbar\xi=(2\pi)^{-k}d\xi$.  Writing
$e^{ix\cdot\xi}=\ang{\xi}^{-2N}(1-\Delta_x)^Ne^{ix\cdot\xi}$ and integrating by
parts shows that the inner integral is $O(\ang\xi^{M-2N})$, so the outer
integral converges for $2N>M+k$ and $\osc{c}\in\CmI(\R^n)$.  When $M<-k$ the
integral in \eqref{conormal:eq:osc} converges absolutely and $\osc{c}$ is
continuous.  By \cref{conormal:lem:density}, $\osc{c_\epsilon}\to\osc{c}$ in
$\CmI$ if $c_\epsilon\to c$ in $S^{M'}$, which justifies the formal
manipulations below.

\begin{lemma}\label{conormal:lem:osc-ops}
  For $c\in S^M(\R^n;\R^k)$:
  \begin{enumerate}[label=(\roman*)]
  \item $D_{y_l}\osc{c}=\osc{D_{y_l}c}$, $D_{x_j}\osc{c}=\osc{\xi_jc+D_{x_j}c}$,
    $x_j\osc{c}=\osc{-D_{\xi_j}c}$ and $\psi\osc{c}=\osc{\psi c}$ for
    $\psi\in\CIb(\R^n)$;
  \item if $M<-k-r$ then $\osc{c}\in\mathcal C^r(\R^n)$ with bounded derivatives
    of order $\le r$;
  \item $\osc{c}$ is smooth in $x\ne0$ and, for every $\epsilon>0$,
    $\sup_{|x|\ge\epsilon}\ang{x}^N|D_x^\alpha D_y^\gamma\osc{c}|<\infty$ for all
    $N,\alpha,\gamma$.
  \end{enumerate}
\end{lemma}
\begin{proof}
  (i) holds for $c$ of order $<-k$ by differentiation under the integral and,
  for $x_j$, by $x_je^{ix\cdot\xi}=D_{\xi_j}e^{ix\cdot\xi}$ and integration by
  parts; it extends by continuity.  (ii) holds since then $D_x^\alpha D_y^\gamma
  \osc{c}$, $|\alpha+\gamma|\le r$, is $\osc{c'}$ with $c'$ of order $<-k$.  (iii):
  by (i), $|x|^{2N}D_x^\alpha D_y^\gamma\osc{c}=\osc{c'}$ with $c'$ of order
  $M+|\alpha|-2N$, which is bounded with all derivatives up to order $r$ once
  $2N>M+|\alpha|+k+r$.
\end{proof}

\begin{lemma}[Reduction]\label{conormal:lem:reduction}
  If $c\in S^M(\R^n;\R^k)$ and $a\in S^M(\R^d;\R^k)$ satisfies
  \begin{equation}\label{conormal:eq:reduction}
    a(y,\xi)\sim\sum_{\alpha}\frac1{\alpha!}(-D_\xi)^\alpha\pa_x^\alpha
    c(0,y,\xi),
  \end{equation}
  then $\osc{c}-\osc{a}\in\CI(\R^n)$.  If $c$ is classical, so is $a$.
\end{lemma}
\begin{proof}
  By Taylor's formula in $x$,
  $c(x,y,\xi)=\sum_{|\alpha|<N}\frac{x^\alpha}{\alpha!}\pa_x^\alpha c(0,y,\xi)
  +\sum_{|\alpha|=N}x^\alpha r_\alpha(x,y,\xi)$ with
  $r_\alpha=\frac N{\alpha!}\int_0^1(1-t)^{N-1}\pa_x^\alpha c(tx,y,\xi)\,dt\in
  S^M(\R^n;\R^k)$.  By \cref{conormal:lem:osc-ops}(i),
  $\osc{x^\alpha f}=\osc{(-D_\xi)^\alpha f}$, so
  \begin{equation*}
    \osc{c}-\osc{a}=\osc{a_N-a}+\sum_{|\alpha|=N}\osc{(-D_\xi)^\alpha r_\alpha},
    \quad a_N=\sum_{|\alpha|<N}\frac1{\alpha!}(-D_\xi)^\alpha\pa_x^\alpha
    c(0,y,\xi).
  \end{equation*}
  Both terms on the right have amplitudes of order $M-N$, so are in
  $\mathcal C^r$ if $N>M+k+r$.  As $N$ is arbitrary, $\osc{c}-\osc{a}$ is smooth.
\end{proof}

\begin{definition}\label{conormal:def:local}
  Let $\Omega\subset\R^n$ be open.  The space $\Ic{m}_{\mathrm c}(\Omega,Y)$
  consists of the $u\in\CmIc(\Omega)$ of the form $u=\osc{c}+f$ with
  $c\in S^M(\R^n;\R^k)$, $M$ given by \eqref{conormal:eq:order}, and
  $f\in\CI(\R^n)$; $\Ic{m}_{\mathrm c,\mathrm{ph}}(\Omega,Y)$ is defined with
  $c\in\Symph{M}(\R^n;\R^k)$.  Further, $\Ic{m}(\Omega,Y)$ is the space of
  $u\in\CmI(\Omega)$ such that $\phi u\in\Ic{m}_{\mathrm c}(\Omega,Y)$ for all
  $\phi\in\CIc(\Omega)$, and similarly in the classical case.
\end{definition}

If $u=\osc{c}+f\in\Ic{m}_{\mathrm c}(\Omega,Y)$ and $\chi\in\CIc(\Omega)$ equals
one near $\supp u$, then $u=\osc{\chi c}+\chi f$, so the amplitude can be taken
with compact support in $(x,y)$.  By \cref{conormal:lem:reduction} it can also be
taken independent of $x$: $u=\osc{a}+f'$ with $a\in S^M_{\mathrm c}(\R^d;\R^k)$
and $f'$ smooth (not compactly supported in general).  We call such an $a$ a
\emph{full symbol} of $u$.

\begin{discussion}
  Two examples with $k=1$ fix the ideas. The delta function $\delta(x_1)$ has
  full symbol $1$, so $M=0$ (see \cref{conormal:ex:delta}). The Heaviside
  function satisfies $H(x_1)-\frac12=\frac12\sgn(x_1)$, and the inverse
  Fourier transform in $x_1$ of $(1-\chi(\xi_1))/(i\xi_1)$, $\chi\in\CIc(\R)$
  equal to one near $0$, differs from $\frac12\sgn(x_1)$ by a smooth function.
  So, after a cut-off, $H(x_1)$ is conormal to $Y=\{x_1=0\}$ with full
  symbol $(1-\chi(\xi_1))/(i\xi_1)$, of order $M=-1$: integrating once lowers
  the order by one. In both cases the full symbol does not depend on $y$; in
  general it may, smoothly.
\end{discussion}

\begin{proposition}\label{conormal:prop:local-basic}
  Let $\Omega\subset\R^n$ be open.  Then (and likewise for the classical spaces):
  \begin{enumerate}[label=(\roman*)]
  \item $\Ic{m}(\Omega,Y)$ is a $\CI(\Omega)$-module, $\Ic{m-1}\subset\Ic{m}$,
    $x_j:\Ic{m}\to\Ic{m-1}$, $D_{x_j}:\Ic{m}\to\Ic{m+1}$, and every smooth vector
    field tangent to $Y$ maps $\Ic{m}$ to itself;
  \item each $u\in\Ic{m}(\Omega,Y)$ is smooth in $\Omega\setminus Y$, and
    $\bigcap_m\Ic{m}(\Omega,Y)=\CI(\Omega)$;
  \item a full symbol of $u\in\Ic{m}_{\mathrm c}(\Omega,Y)$ is determined modulo
    $S^{-\infty}$; and $u\in\Ic{m-1}_{\mathrm c}(\Omega,Y)$ if and only if its
    full symbols lie in $S^{M-1}$;
  \item if $u_j\in\Ic{m-j}_{\mathrm c}(\Omega,Y)$ have supports in a fixed
    compact set, there is $u\in\Ic{m}_{\mathrm c}(\Omega,Y)$ with
    $u-\sum_{j<N}u_j\in\Ic{m-N}_{\mathrm c}(\Omega,Y)$ for all $N$;
  \item if $P=\sum_{|\alpha|\le q}p_\alpha(z)D_z^\alpha$ with
    $p_\alpha\in\CI(\Omega)$, then $P:\Ic{m}(\Omega,Y)\to\Ic{m+q}(\Omega,Y)$ and,
    if $u\in\Ic{m}_{\mathrm c}$ has full symbol $a$, then $Pu$ has full symbol
    congruent to $p_q(0,y;\xi,0)a(y,\xi)$ modulo $S^{M+q-1}$, where
    $p_q(z,\zeta)=\sum_{|\alpha|=q}p_\alpha(z)\zeta^\alpha$.
  \end{enumerate}
\end{proposition}
\begin{proof}
  (i): for $\psi\in\CI(\Omega)$ and $u\in\Ic{m}_{\mathrm c}$,
  $\psi u=\psi\chi u=\osc{\psi\chi c}+\psi\chi f$ with $\chi$ as above.  The rest
  follows from \cref{conormal:lem:osc-ops}(i), since the vector fields tangent
  to $Y$ are spanned over $\CI$ by $x_i\pa_{x_j}$ and $\pa_{y_l}$ (see
  \cref{conormal:lem:generators} below) and $x_i\pa_{x_j}$ maps an amplitude of
  order $M$ to $(-D_{\xi_i})(i\xi_jc+\pa_{x_j}c)$, of order $M$.
  (ii): smoothness off $Y$ is \cref{conormal:lem:osc-ops}(iii).  If $u$ lies in
  every $\Ic{m}$, then for each $\phi$ and $r$, $\phi u$ is the sum of a smooth
  function and $\osc{c}$ with $c$ of order $<-k-r$, so $\phi u\in\mathcal C^r$.
  Conversely $\CIc(\Omega)\subset\Ic{m}_{\mathrm c}$ trivially.
  (iii): suppose $\osc{a}\in\CI(\R^n)$ with $a\in S^M_{\mathrm c}(\R^d;\R^k)$.
  Let $\chi\in\CIc(\R^k)$ equal one near $0$.  For each $y$, $\osc{a}(\cdot,y)$
  is the inverse Fourier transform of $a(y,\cdot)$, and it is the sum of
  $\chi(x)\osc{a}(x,y)$, smooth with compact support, and
  $(1-\chi(x))\osc{a}(x,y)$, which is in $\cS(\R^k)$ by
  \cref{conormal:lem:osc-ops}(iii), both with seminorms bounded with all
  $y$-derivatives.  Hence $a\in S^{-\infty}$.  If now $u=\osc{a}+f=\osc{b}+g$
  then $\osc{a-b}$ is smooth, so $a-b\in S^{-\infty}$.  The second statement
  follows: if $u\in\Ic{m-1}_{\mathrm c}$ it has a full symbol $b\in S^{M-1}$, and
  any other full symbol differs from $b$ by an element of $S^{-\infty}$.
  (iv): with full symbols $a_j\in S^{M-j}_{\mathrm c}$ supported in a fixed
  compact set of $y$, take $u=\chi\osc{a}$ with $a\sim\sum a_j$ (with support in
  the same set) and $\chi\in\CIc(\Omega)$ equal one near all the supports.
  (v): by (i) and \cref{conormal:lem:osc-ops}(i), $D^\alpha\osc{a}=\osc{(\xi,0)^\alpha
  a}+\osc{c'}$ with $c'$ of order $M+|\alpha|-1$ (the $y$-derivatives do not
  raise the order), and $p_\alpha\osc{b}=\osc{p_\alpha b}$, whose reduction is
  $p_\alpha(0,y)b$ modulo order $M'-1$ by \eqref{conormal:eq:reduction}.
\end{proof}

\begin{remark}\label{conormal:rem:point-local}
  For $k=n$ (so $Y=\{0\}$) we recover the spaces of \cref{conormal:sec:point}:
  $\Ic{m}_\cS(\R^n,\{0\})=\Ic{m}_{\mathrm c}(\R^n,\{0\})+\cS(\R^n)$.  Indeed if
  $\hat u\in S^{m-\frac n4}$ then $u=\chi u+(1-\chi)u$ with $(1-\chi)u\in\cS$ by
  \cref{conormal:lem:point-basic}(iii) and $\chi u=\osc{\chi\hat u}$.  Conversely
  $\osc{c}$ with $c$ compactly supported in $z$ has Fourier transform
  $\int\hat c(\zeta-\xi;\xi)\,\dbar\xi$, where $\hat c(\cdot;\xi)$ is the Fourier
  transform of $c(\cdot,\xi)$; this is in $S^{m-\frac n4}$ by Peetre's inequality
  $\ang{\xi}^{M}\le2^{|M|}\ang{\zeta}^M\ang{\zeta-\xi}^{|M|}$ (moving derivatives
  onto $c$ after the substitution $\xi\mapsto\zeta-\xi$).  In particular
  $\Ic{m}_\cS(\R^n,\{0\})$ is a $\CIb(\R^n)$-module, and in particular an
  $\cS(\R^n)$-module.
\end{remark}

\begin{example}[An elliptic problem at a point]\label{conormal:ex:elliptic-point}
  Let $P=\sum_{|\alpha|\le q}p_\alpha(z)D^\alpha$ with $p_\alpha\in\CIb(\R^n)$,
  and suppose $P$ is \emph{elliptic at $0$}: $p_q(0,\zeta)=\sum_{|\alpha|=q}
  p_\alpha(0)\zeta^\alpha\ne0$ for $\zeta\ne0$.  Then for every
  $f\in\Ic{m}_\cS(\R^n,\{0\})$ there is $u\in\Ic{m-q}_\cS(\R^n,\{0\})$ with
  $Pu-f\in\cS(\R^n)$.  Indeed, by \cref{conormal:prop:local-basic}(v) and
  \cref{conormal:rem:point-local}, $P:\Ic{r}_\cS\to\Ic{r+q}_\cS$ and
  $Pv-p_q(0,D)v\in\Ic{r+q-1}_\cS$ for $v\in\Ic{r}_\cS$.  Choose
  $\chi\in\CIc(\R^n)$ equal to one near $0$ and set
  $\hat u_0=(1-\chi)\hat f/p_q(0,\cdot)\in S^{m-q-\frac n4}$, using that
  $(1-\chi)/p_q(0,\cdot)\in\Symph{-q}$.  Then $p_q(0,D)u_0-f\in\cS$, so
  $Pu_0=f-f_1$ with $f_1\in\Ic{m-1}_\cS$.  Repeating with $f$ replaced by $f_1$,
  we find $u_j\in\Ic{m-q-j}_\cS$ with $P(u_0+\dots+u_j)-f\in\Ic{m-j-1}_\cS$;
  an asymptotic sum $u\sim\sum u_j$ (\cref{conormal:prop:local-basic}(iv)) then
  satisfies $Pu-f\in\bigcap_j\Ic{m-j}_\cS=\cS(\R^n)$.  The same works in the
  classical spaces.
\end{example}

\subsection{Pseudodifferential operators}
We now let pseudodifferential operators on $\R^n$, as in \cref{ch:euclidean},
act on conormal distributions.  For $a\in S^M_{\mathrm c}(\R^d;\R^k)$ the
distribution $\osc{a}$ is tempered, with Fourier transform
\begin{equation}\label{conormal:eq:FT-osc}
  \widehat{\osc{a}}(\xi,\eta)=\tilde a(\xi,\eta)=\int e^{-iy\cdot\eta}
  a(y,\xi)\,dy,\qquad
  |D_\xi^\beta D_\eta^\gamma\tilde a(\xi,\eta)|\le
  C_{N\beta\gamma}\ang{\eta}^{-N}\ang\xi^{M-|\beta|},
\end{equation}
the estimate following by integration by parts in $y$.

\begin{proposition}\label{conormal:prop:psido-local}
  Let $A=\Op(b)$, $b\in S^{m'}(\R^n;\R^n)$, and $u\in\Ic{m}_{\mathrm c}(\R^n,Y)$
  with full symbol $a\in S^M_{\mathrm c}(\R^d;\R^k)$.  Then
  $Au\in\Ic{m+m'}(\R^n,Y)$; more precisely $Au-\osc{c}\in\cS(\R^n)$ where
  $c\in S^{M+m'}(\R^n;\R^k)$ is given by
  \begin{equation}\label{conormal:eq:psido-amp}
    c(x,y,\xi)=\int e^{iy\cdot\eta}b(x,y,\xi,\eta)\tilde a(\xi,\eta)\,\dbar\eta
    \sim\sum_\gamma\frac1{\gamma!}\pa_\eta^\gamma b(x,y,\xi,0)\,D_y^\gamma
    a(y,\xi).
  \end{equation}
  In particular $Au$ has full symbol (if $b$ has compact support in $z$)
  \begin{equation}\label{conormal:eq:psido-full}
    a'(y,\xi)\sim\sum_{\alpha,\gamma}\frac1{\alpha!\gamma!}(-D_\xi)^\alpha\pa_x^\alpha
    \Big(\pa_\eta^\gamma b(x,y,\xi,0)D_y^\gamma a(y,\xi)\Big)\Big|_{x=0},
  \end{equation}
  and $a'-b(0,y,\xi,0)a\in S^{M+m'-1}$.  If $b$ and $a$ are classical, so is $c$.
\end{proposition}
\begin{proof}
  Write $u=\osc{a}-g$.  Then $g$ is smooth, and outside a compact set equals
  $\osc{a}$, which vanishes for $y$ outside the support of $a$ and is rapidly
  decreasing in $x$ by \cref{conormal:lem:osc-ops}(iii); so $g\in\cS(\R^n)$
  and $Ag\in\cS(\R^n)$.  By \eqref{conormal:eq:FT-osc}, $A\osc{a}(x,y)=\int
  e^{ix\cdot\xi+iy\cdot\eta}b(x,y,\xi,\eta)\tilde a(\xi,\eta)\,\dbar\xi\,\dbar\eta$,
  the $\eta$-integral being absolutely convergent; this is $\osc{c}$ with $c$ as in
  \eqref{conormal:eq:psido-amp}.  To estimate $c$, note that derivatives in
  $(x,y)$ and $\xi$ may be taken under the integral, and by Peetre's inequality
  \begin{equation*}
    |\pa_z^\delta\pa_\xi^{\beta}b(z,\xi,\eta)|\le
    C\ang{(\xi,\eta)}^{m'-|\beta|}\le C'\ang\xi^{m'-|\beta|}
    \ang{\eta}^{|m'|+|\beta|};
  \end{equation*}
  combined with \eqref{conormal:eq:FT-osc} (and $y$-derivatives falling on
  $e^{iy\cdot\eta}$ producing powers of $\eta$) this gives $c\in S^{M+m'}$.  For
  the expansion, Taylor's formula in $\eta$ gives
  $b(z,\xi,\eta)=\sum_{|\gamma|<L}\frac{\eta^\gamma}{\gamma!}\pa_\eta^\gamma
  b(z,\xi,0)+\sum_{|\gamma|=L}\eta^\gamma r_\gamma(z,\xi,\eta)$ with
  \begin{equation*}
    |\pa_z^\delta\pa_\xi^\beta r_\gamma|\le C\sup_{0\le t\le1}
    \ang{(\xi,t\eta)}^{m'-L-|\beta|}\le C\ang\xi^{m'-L-|\beta|}
    \ang\eta^{|m'-L|+|\beta|}.
  \end{equation*}  Since $\int e^{iy\cdot\eta}\eta^\gamma\tilde
  a\,\dbar\eta=D_y^\gamma a$, the terms with $|\gamma|<L$ give the expansion
  and the remainder is of order $M+m'-L$ by the same estimate as before.
  Finally \eqref{conormal:eq:psido-full} follows from
  \cref{conormal:lem:reduction}.  In the classical case each term of
  \eqref{conormal:eq:psido-amp} is classical, so the asymptotic sum is.
\end{proof}

The microlocal content of this computation is that $u$ is singular only in the
conormal directions to $Y$.  Recall from \cref{ch:microlocal} that
$(\bar z,\bar\zeta)\notin\WF(u)$ if and only if there is $A\in\Psido{0}(\R^n)$
elliptic at $(\bar z,\bar\zeta)$ with $Au\in\CI$.  The conormal bundle of $Y$
is $N^*Y=\{(0,y;\xi,0)\}\subset T^*\R^n$.

\begin{figure}[htbp]
\centering
\begin{tikzpicture}[>=Latex,scale=0.9]
  \begin{scope}
    \draw[->] (-2.3,0)--(2.5,0) node[right]{$y$};
    \draw[->] (0,-2)--(0,2.2) node[above]{$x$};
    \draw[very thick] (-2.2,0)--(2.2,0) node[below left,font=\footnotesize]{$Y$};
    \foreach \y in {-1.6,-0.8,0.8,1.6}{
      \foreach \x in {0.5,1,1.5}{
        \draw[->] (\y,\x)--(\y,\x+0.3*\x);
        \draw[->] (\y,-\x)--(\y,-\x-0.3*\x);}}
    \foreach \x in {-1.5,-0.5,0.5,1.5} \draw[->,gray] (-2.1,\x)--(-1.7,\x);
    \node[font=\footnotesize] at (0,-2.6) {(a) $x\pa_x$ (black) and $\pa_y$ (gray)};
  \end{scope}
  \begin{scope}[xshift=6.5cm]
    \draw[->] (-2.3,0)--(2.5,0) node[right]{$y$};
    \draw[->] (0,-2)--(0,2.2) node[above]{$x$};
    \draw[very thick] (-2.2,0)--(2.2,0) node[below left,font=\footnotesize]{$Y$};
    \foreach \y in {-1.6,-0.8,0.8,1.6} \draw[<->,thick] (\y,-0.8)--(\y,0.8);
    \node[font=\footnotesize] at (0,-2.6) {(b) $N^*Y=\{(0,y;\xi,0)\}$};
  \end{scope}
\end{tikzpicture}
\caption[Vector fields tangent to $Y$ and the conormal bundle]{The local model
with $k=d=1$. (a) Vector fields tangent to $Y=\{x=0\}$ are generated by
$x_i\pa_{x_j}$ and $\pa_{y_l}$ (\cref{conormal:lem:generators}); they preserve
each $\Ic{m}(\Omega,Y)$ (\cref{conormal:prop:local-basic}(i)), and repeated
application of them characterizes conormality (\cref{conormal:sec:iterated}).
(b) The conormal bundle, which contains the wavefront set of every conormal
distribution (\cref{conormal:prop:WF-local}).}
\label{conormal:fig:local}
\end{figure}

\begin{proposition}\label{conormal:prop:WF-local}
  If $u\in\Ic{m}(\Omega,Y)$ then $\WF(u)\subset N^*Y\setminus0$.  More
  precisely, if $u\in\Ic{m}_{\mathrm c}(\Omega,Y)$ has full symbol $a$, then
  $(0,\bar y;\bar\xi,0)\notin\WF(u)$ if and only if $a$ satisfies the estimates
  of $S^{-\infty}$ in $U\times\Gamma$, for some neighbourhood $U$ of $\bar y$ and
  open cone $\Gamma\ni\bar\xi$.
\end{proposition}
\begin{proof}
  We may assume $u\in\Ic{m}_{\mathrm c}$.  Over $\Omega\setminus Y$, $u$ is
  smooth.  Let $\bar z=(0,\bar y)$ and $\bar\zeta=(\bar\xi,\bar\eta)$ with
  $\bar\eta\ne0$.  Choose $\delta>0$ with $|\bar\eta|>2\delta|\bar\xi|$, and
  $b=\phi(z)\psi(\zeta)$ with $\phi\in\CIc$, $\phi(\bar z)\ne0$,
  $\psi\in\Symph{0}(\R^n)$ homogeneous of degree $0$ in $|\zeta|\ge1$, equal to
  one near $\bar\zeta/|\bar\zeta|$ there, and supported in
  $\{|\eta|\ge\delta|\xi|,\ |\zeta|\ge\frac12\}$.  On the support of $\psi$,
  $\ang\eta\ge c\ang\xi$, so by \eqref{conormal:eq:FT-osc} the amplitude $c$ of
  \eqref{conormal:eq:psido-amp} satisfies
  $|c|\le C_N\int_{\ang\eta\ge c\ang\xi}\ang\eta^{-N}\ang\xi^{M}\,d\eta\le
  C'_N\ang\xi^{M-N+d}$ for all $N$, and similarly for derivatives; thus $c\in
  S^{-\infty}$ and $\Op(b)u\in\CI$.  Hence $\WF(u)\subset N^*Y$.

  For the refined statement let $\bar\zeta=(\bar\xi,0)$.  If $a$ is of order
  $-\infty$ in $U\times\Gamma$, take $b=\phi(z)\psi(\xi,\eta)$ as before but
  with $\phi$ supported in $\R^k\times U'$, $U'\Subset U$, and $\psi$ supported in
  a small conic neighbourhood of $(\bar\xi,0)$, so that $(y,\xi)\in U'\times\Gamma$
  on the support of $\phi(x,y)\psi(\xi,0)$ for $|\xi|\ge1$.  Each term in
  \eqref{conormal:eq:psido-full} then vanishes to infinite order (the terms involve
  $a$ only at points of $U'\times\Gamma$), so $Au\in\CI$ and $A$ is elliptic at
  $(\bar z,\bar\zeta)$.  Conversely, if $(\bar z,\bar\zeta)\notin\WF(u)$, choose
  $B=\Op(b)$ as just described with $b=1$ on a conic neighbourhood of
  $(\bar z,\bar\zeta)$ and $\WFp(B)\cap\WF(u)=\emptyset$; then $Bu\in\CI$ by
  \cref{ch:microlocal}, so its full symbol \eqref{conormal:eq:psido-full} is in
  $S^{-\infty}$, and near $(\bar y,\bar\xi)$ all terms with
  $(\alpha,\gamma)\neq0$ vanish because $b$ is constant there, leaving
  $a\in S^{-\infty}$ near $(\bar y,\bar\xi)$.
\end{proof}

\section{Invariance}\label{conormal:sec:invariance}

The definition of the local model uses linear coordinates and a Fourier
transform, so it is not clear that it makes sense on a manifold. We now show that the spaces $\Ic{m}_{\mathrm c}(\Omega,Y)$ are invariant under
diffeomorphisms preserving $Y$, and compute the effect on the full symbol
modulo lower order terms.  Let $F:\Omega'\to\Omega$ be a diffeomorphism between
open subsets of $\R^n$ with $F(\Omega'\cap Y)=\Omega\cap Y$.  Write
$F(x,y)=(F'(x,y),F''(x,y))$; then $F'(0,y)=0$, so $\pa_yF'(0,y)=0$, and
\begin{equation}\label{conormal:eq:F-data}
  g(y)=F''(0,y),\qquad A(y)=\pa_xF'(0,y)\in\GL(k,\R)
\end{equation}
define a diffeomorphism $g$ of $\Omega'\cap Y$ onto $\Omega\cap Y$ and the
linear maps $A(y)$ induced by $F_*$ on the normal spaces.

The proof of invariance splits $F$ into two pieces of very different
nature:
\begin{equation*}
\begin{tikzcd}[column sep=small,row sep=normal,cells={nodes={font=\small}}]
  & F=F_0\circ H\arrow[dl]\arrow[dr] & \\
  \begin{array}{c}F_0(x,y)=(A(y)x,g(y))\\ \text{fibre-linear: exact formula}\\
  \text{(\cref{conormal:lem:fibre-linear})}\end{array}\arrow[dr]
  & & \begin{array}{c}H=\Id+O(\abs x^2)\text{ in }x,\ O(\abs x)\text{ in }y\\
  \text{Moser homotopy: }H^*u-u\in\Ic{m-1}\\
  \text{(\cref{conormal:prop:moser})}\end{array}\arrow[dl]\\
  & \text{\cref{conormal:thm:invariance}: }F^*u\in\Ic{m},\ 
  \text{symbol }\eqref{conormal:eq:symbol-transform} &
\end{tikzcd}
\end{equation*}
The fibre-linear part acts on the symbol by a linear change of variables, and
the remaining part, which is the identity to first order on $Y$, only changes
lower order terms.

\begin{lemma}\label{conormal:lem:fibre-linear}
  Let $V\subset\R^d$ be open, let $g$ be a diffeomorphism of $V$ onto $g(V)$
  and let $A\in\CI(V;\GL(k,\R))$.  Set $F_0(x,y)=(A(y)x,g(y))$, a diffeomorphism
  of $\R^k\times V$ onto $\R^k\times g(V)$.  If $a\in S^M_{\mathrm c}(\R^d;\R^k)$ has
  support in $g(V)\times\R^k$ then
  \begin{equation*}
    F_0^*\osc{a}=\osc{a_{F}},\qquad
    a_{F}(y,\eta)=|\det A(y)|^{-1}a\big(g(y),A(y)^{-t}\eta\big)
    \in S^M_{\mathrm c}(\R^d;\R^k),
  \end{equation*}
  where $A^{-t}=(A^t)^{-1}$; $a_F$ is classical if $a$ is.
\end{lemma}
\begin{proof}
  For fixed $y$ this is the linear change of variables
  $\xi=A(y)^{-t}\eta$ in $\int e^{iA(y)x\cdot\xi}a(g(y),\xi)\,\dbar\xi$.  That
  $a_F$ is a symbol follows from the invariance of symbol estimates under
  smooth families of linear maps, the support of $a_F$ in $y$ being compact.
\end{proof}

\begin{lemma}\label{conormal:lem:tangent-lower}
  A vector field of the form
  \begin{equation}\label{conormal:eq:V-form}
    V=\sum_{i,j,l}x_ix_j\alpha_{ijl}(z)\pa_{x_l}+\sum_{i,l}x_i\beta_{il}(z)\pa_{y_l},
    \qquad\alpha_{ijl},\beta_{il}\in\CI(\Omega),
  \end{equation}
  maps $\Ic{m}(\Omega,Y)$ to $\Ic{m-1}(\Omega,Y)$.
\end{lemma}
\begin{proof}
  By \cref{conormal:prop:local-basic}(i), $x_ix_j\pa_{x_l}$ maps $\Ic{m}$ into
  $\Ic{m+1-2}$ and $x_i\pa_{y_l}$ maps $\Ic{m}$ into $\Ic{m-1}$, and the
  coefficients are harmless.
\end{proof}

We also need conormal distributions depending smoothly on a parameter
$s\in[0,1]$.  A \emph{smooth family} $u(s)\in\Ic{m}_{\mathrm c}(\Omega,Y)$ is
one of the form $u(s)=\osc{c(s)}+f(s)$ with $c\in S^M(\R_s\times\R^n;\R^k)$
and $f\in\CI(\R\times\R^n)$, the supports of $u(s)$ lying in a fixed compact
subset of $\Omega$; that is, $u$ is an element of $\Ic{m}_{\mathrm c}$ on
$\R_s\times\R^n$ with respect to $\R_s\times Y$, restricted to $[0,1]$.  By
\cref{conormal:prop:local-basic} applied in $n+1$ variables, $\pa_su(s)$ and
$\int_s^1u(s')\,ds'$ are smooth families of order $m$, $V_su(s)$ is a smooth
family of order $m-1$ if $V_s$ is of the form \eqref{conormal:eq:V-form} with
coefficients smooth in $s$, and families of orders $m-j$ have asymptotic sums.

\begin{proposition}\label{conormal:prop:moser}
  Let $H$ be a diffeomorphism from a neighbourhood $O$ of $\bar z=(0,\bar y)$
  onto its image such that $H(0,y)=(0,y)$ and $\pa_xH'(0,y)=\Id$ for
  $(0,y)\in O$.  Then there is a neighbourhood $O'$ of $\bar z$ such that
  $H^*u-u\in\Ic{m-1}_{\mathrm c}(\R^n,Y)$ for all
  $u\in\Ic{m}_{\mathrm c}(\R^n,Y)$ with support in $O'$; similarly in the
  classical case.
\end{proposition}
\begin{discussion}
  The method is Moser's homotopy trick. Join $H$ to the identity by the family
  $H_s$. The derivative $\frac d{ds}H_s^*$ is $H_s^*$ composed with a vector field
  $V_s$, and because $H$ agrees with the identity to first order on $Y$, $V_s$
  has the special form \eqref{conormal:eq:V-form}, which lowers the conormal
  order by one (\cref{conormal:lem:tangent-lower}). Instead of integrating
  $\frac d{ds}H_s^*u$ directly, we solve the transport equation
  $\pa_su(s)+V_su(s)\approx0$ with $u(1)=u$; then $H_s^*u(s)$ is almost constant
  in $s$, and $H^*u\approx u(0)$, which differs from $u$ by terms of lower
  order.
\end{discussion}

\begin{proof}
  Set $H_s(z)=z+s(H(z)-z)$, $s\in[0,1]$, so $H_0=\Id$, $H_1=H$ and
  $H_s(0,y)=(0,y)$.  At points of $Y$, $dH_s=\begin{psmallmatrix}\Id&0\\
  sB&\Id\end{psmallmatrix}$ with $B=\pa_xH''(0,y)$, which is invertible.  By the
  inverse function theorem, applied uniformly for $s$ in the compact set
  $[0,1]$, there is a ball $O_1\ni\bar z$ such that each $H_s$ is a
  diffeomorphism of $O_1$ onto its image, and a smaller ball $O'$ contained in
  $H_s(O_1)$ for all $s$.  For $w\in O'$ let
  \begin{equation*}
    V_s(w)=(H-\Id)\big(H_s^{-1}(w)\big),
  \end{equation*}
  a vector field depending smoothly on $s$.  For smooth $v$ the chain rule gives
  \begin{equation}\label{conormal:eq:chain}
    \frac{d}{ds}H_s^*v=H_s^*(V_sv),
  \end{equation}
  and the same holds for a smooth family of distributions $v(s)$ in the form
  $\frac d{ds}H_s^*v(s)=H_s^*(\pa_sv+V_sv)$: pairing with a test function
  $\phi$, $\pair{H_s^*v(s)}\phi=\pair{v(s)}{\phi_s}$ with
  $\phi_s=|\det dH_s^{-1}|\,(H_s^{-1})^*\phi$ smooth in $s$, and the identity
  $\pair{w}{\pa_s\phi_s}=\pair{H_s^*(V_sw)}{\phi}$, valid for smooth $w$ by
  \eqref{conormal:eq:chain}, extends to distributions by density.

  Since $H'(x,y)-x=O(|x|^2)$, $H''(x,y)-y=O(|x|)$ and the $x$-component of
  $H_s^{-1}(w)$ vanishes on $Y$ (so is a combination $\sum x_i(\dots)$), $V_s$
  is of the form \eqref{conormal:eq:V-form} with coefficients smooth in $s$.

  Now let $u\in\Ic{m}_{\mathrm c}$ have support in $O'$.  We construct a smooth
  family $u(s)$, with supports in a fixed compact subset of $O'$, such that
  \begin{equation}\label{conormal:eq:moser}
    u(1)=u,\qquad \pa_su(s)+V_su(s)=r(s)\in\CIc([0,1]\times O'),\qquad
    u(0)-u\in\Ic{m-1}_{\mathrm c}.
  \end{equation}
  Set $u_0(s)=u$ and $f_1(s)=V_su\in\Ic{m-1}$.  Inductively, given smooth
  families $u_0,\dots,u_N$ with $u_j$ of order $m-j$, $u_j(1)=0$ for $j\ge1$,
  and $\pa_s(u_0+\dots+u_N)+V_s(u_0+\dots+u_N)=f_{N+1}$ of order $m-N-1$, put
  $u_{N+1}(s)=\int_s^1f_{N+1}(s')\,ds'$.  Then $u_{N+1}$ has order $m-N-1$,
  vanishes at $s=1$, and the new error is $f_{N+2}=V_su_{N+1}$, of order
  $m-N-2$.  Let $\tilde u(s)\sim\sum_ju_j(s)$ be an asymptotic sum, multiplied by a
  cut-off equal to one near the supports.  Then $\pa_s\tilde u+V_s\tilde u$ is a
  smooth family of every order, hence a smooth function of $(s,z)$, and
  $\tilde u(1)-u=\tilde u(1)-\sum_{j<N}u_j(1)$ has order $m-N$ for every $N$, so
  is smooth.  Thus $u(s)=\tilde u(s)-(\tilde u(1)-u)$ satisfies
  \eqref{conormal:eq:moser}; note that $u(0)-u\equiv\sum_{j\ge1}u_j(0)$ modulo
  $\CI$ has order $m-1$.
  Then
  \begin{equation*}
    H^*u=H_1^*u(1)=H_0^*u(0)+\int_0^1\frac{d}{ds}H_s^*u(s)\,ds
    =u(0)+\int_0^1H_s^*r(s)\,ds,
  \end{equation*}
  and the last integral is smooth.  Hence
  $H^*u-u=(u(0)-u)+\CI\subset\Ic{m-1}_{\mathrm c}$.
\end{proof}

\begin{theorem}\label{conormal:thm:invariance}
  Let $F:\Omega'\to\Omega$ be a diffeomorphism of open subsets of $\R^n$ with
  $F(\Omega'\cap Y)=\Omega\cap Y$, and $g$, $A$ as in \eqref{conormal:eq:F-data}.
  Then $F^*:\Ic{m}_{\mathrm c}(\Omega,Y)\to\Ic{m}_{\mathrm c}(\Omega',Y)$, and if
  $u\in\Ic{m}_{\mathrm c}(\Omega,Y)$ has full symbol $a$, then $F^*u$ has full
  symbol congruent modulo $S^{M-1}$ to
  \begin{equation}\label{conormal:eq:symbol-transform}
    |\det A(y)|^{-1}a\big(g(y),A(y)^{-t}\eta\big).
  \end{equation}
  The same holds for the classical spaces.
\end{theorem}
\begin{proof}
  Using a partition of unity and the module property we may assume that $u$ has
  support in a small neighbourhood of $F(\bar z)$, $\bar z\in Y\cap\Omega'$ (away
  from $Y$, $u$ and $F^*u$ are smooth).  Let $F_0(x,y)=(A(y)x,g(y))$, a
  diffeomorphism of $\R^k\times(Y\cap\Omega')$ onto $\R^k\times(Y\cap\Omega)$
  (identifying $Y$ with $\R^d$), and $H=F_0^{-1}\circ F$, defined near $\bar z$.
  Then $H(0,y)=(0,y)$ and $\pa_xH'(0,y)=A(y)^{-1}\pa_xF'(0,y)=\Id$, so
  \cref{conormal:prop:moser} applies to $H$.  With $u=\osc{a}+f$ we have, by
  \cref{conormal:lem:fibre-linear}, $F_0^*u=\osc{a_F}+F_0^*f$ (after cutting
  off), and $F^*u=H^*(F_0^*u)\in F_0^*u+\Ic{m-1}_{\mathrm c}$.
\end{proof}

\section{Conormality at the zero section of a vector bundle}\label{conormal:sec:bundle}

The zero section of a vector bundle is the intermediate case between the local
model and a general submanifold: the fibres carry a linear structure, so the
Fourier transform in the fibres makes sense, but there are no global
coordinates. By the tubular neighbourhood theorem every submanifold looks like
this near $Y$ (\cref{conormal:sec:manifold}).

Let $\pi:W\to Y$ be a real vector bundle of rank $k$ over a manifold $Y$ of
dimension $d$, so $W$ has dimension $n=d+k$, and let $0_W\subset W$ be the zero
section and $W'$ the dual bundle.  The linear invariance of the radial
compactification allows it to be performed fibrewise: $\RC W$ is a manifold with
boundary, fibred over $Y$ with fibres $\RC{W_y}$ (formally, the bundle associated
to the frame bundle of $W$ with fibre $\RC{\R^k}$), and $\rho^{-M}\CI(\RC W)$
is well defined.  Symbols with bounds $S^M(W)$ on a bundle over a compact base,
or with compact support in the base, are defined via local trivializations, by
the linear invariance of the estimates.

Let $\Omega_{\mathrm{fib}}=\Omega_{\mathrm{fib}}(W')$ be the line bundle over
$W'$ whose fibre at $\zeta\in W'_y$ is $\Omega(W'_y)$, the densities on the fibre.
The fibrewise version of \eqref{conormal:eq:invariant-FT} is an isomorphism
\begin{equation*}
  \Ffib:\cS(W)\longrightarrow\cS(W';\Omega_{\mathrm{fib}}),
\end{equation*}
defined for (fibrewise) Schwartz functions; in a local trivialization
$W|_U\cong U\times\R^k$ it is the partial Fourier transform in the fibre
variables followed by multiplication by $|d\xi|$, and it does not depend on the
trivialization.  Its inverse is defined on symbols as an oscillatory integral.

\begin{definition}\label{conormal:def:bundle}
  For $Y$ compact, $\Ic{m}(W,0_W)$ is the space of distributions $u=\Ffib^{-1}b+f$
  with $b\in S^{M}(W';\Omega_{\mathrm{fib}})$ (that is, $b=\beta|d\xi|$ with
  $\beta\in S^M$ in local trivializations), $M=m+\frac n4-\frac k2$, and
  $f\in\CI(W)$; similarly $\Ic{m}_{\mathrm{ph}}(W,0_W)$ with $b$ classical.
\end{definition}

In a local trivialization, with fibre coordinates $x$ and base coordinates $y$,
$\Ffib^{-1}(\beta|d\xi|)=\osc{\beta}$, so $\Ic{m}(W,0_W)$ consists of the
distributions which are smooth outside $0_W$ and lie in the spaces of
\cref{conormal:def:local} in every such local trivialization.  Changes of
trivialization are of the fibre-linear form of \cref{conormal:lem:fibre-linear}
(with $g$ a change of coordinates on $Y$), under which $\beta|d\xi|$ transforms
as a fibre density, so we obtain a well-defined symbol map with values in
$S^M(W';\Omega_{\mathrm{fib}})/S^{M-1}(W';\Omega_{\mathrm{fib}})$, and, by
\cref{conormal:prop:local-basic}(iii) and the surjectivity of the inverse
Fourier transform, a short exact sequence
\begin{equation}\label{conormal:eq:bundle-exact}
  \Ic{m-1}(W,0_W)\hookrightarrow\Ic{m}(W,0_W)\xrightarrow{\ \sigma_m\ }
  S^{M}(W';\Omega_{\mathrm{fib}})/S^{M-1}(W';\Omega_{\mathrm{fib}}),
\end{equation}
and in the classical case the corresponding sequence onto the smooth fibre
densities on $W'\setminus0$ homogeneous of degree $M+k$.  For a product
$W=Y\times\R^k$ the definition says that $\Ic{m}(Y\times\R^k,Y\times\{0\})$
consists of smooth functions of $y\in Y$ with values in
$\Ic{m-\frac d4}(\R^k,\{0\})$ (up to smooth terms), the topology of
$\CI(Y;\rho^{-M}\CI(\RC{\R^k}))=\rho^{-M}\CI(Y\times\RC{\R^k})$ making precise
what it means for a conormal distribution to depend smoothly on parameters.

Bundle isomorphisms, possibly covering a diffeomorphism of the base, act on
$\Ic{m}(W,0_W)$ by \cref{conormal:lem:fibre-linear}.  For the extension to
submanifolds, the essential point is invariance under diffeomorphisms which do
not preserve the linear structure.  If $h$ is a diffeomorphism between
neighbourhoods of $0_W$ which maps $0_W$ to itself, then at each point of $0_W$
the differential $h_*$ preserves $T0_W$ and so induces a map on
$T_yW/T_y0_W=W_y$.

\begin{lemma}\label{conormal:lem:factor}
  Any diffeomorphism $h$ of a neighbourhood of $0_W$ onto a neighbourhood of
  $0_W$ with $h(0_W)=0_W$ factorizes, near $0_W$, as $h=h_0\circ h_1$ where $h_0$
  is a bundle isomorphism (covering $h|_{0_W}$) and $h_1$ is the identity on
  $0_W$ with $(h_1)_*=\Id$ on $W_y=T_yW/T_y0_W$ for all $y$.
\end{lemma}
\begin{proof}
  Take $h_0$ to be the bundle map induced by $h$ on the normal spaces, i.e.\
  $h_0(w)=h_*(w)$ for $w\in W_y=N_y0_W$, covering $h|_{0_W}$; then
  $h_1=h_0^{-1}\circ h$ has the stated properties.  In local trivializations this
  is the factorization $F=F_0\circ H$ used in the proof of
  \cref{conormal:thm:invariance}.
\end{proof}

\begin{proposition}\label{conormal:prop:bundle-invariance}
  Let $Y$ be compact and $h:N\to N'$ a diffeomorphism between open
  neighbourhoods of $0_W$ which is the identity on $0_W$ and induces the
  identity on $W_y$ for all $y\in Y$.  If $u\in\Ic{m}(W,0_W)$ has support in $N'$
  then $h^*u-u\in\Ic{m-1}(W,0_W)$; in particular $\sigma_m(h^*u)=\sigma_m(u)$.
  Combined with \cref{conormal:lem:factor}, every diffeomorphism $h$ between
  neighbourhoods of $0_W$ mapping $0_W$ to itself acts on conormal distributions
  supported in its range, with $\sigma_m(h^*u)=h_0^*\sigma_m(u)$.
\end{proposition}
\begin{proof}
  Cover $Y$ by finitely many open sets over which $W$ is trivial, take a
  partition of unity $\phi_j$ of a neighbourhood of $0_W$ subordinate to their
  preimages and apply \cref{conormal:prop:moser} (in the form of the proof of
  \cref{conormal:thm:invariance}) to each $\phi_ju$, noting that $u-\sum\phi_ju$
  is smooth.
\end{proof}

\section{Conormal distributions on manifolds}\label{conormal:sec:manifold}

\subsection{Preliminaries}
Manifolds are discussed systematically in \cref{ch:manifolds}; here we need
only a few notions.  Let $X$ be a smooth manifold of dimension $n$.  The
density bundles $\Omega^t(X)$, $t\in\R$, are the line bundles with fibres
$\Omega^t(T_pX)$; in local coordinates their sections are $v|dz|^t$, and under a
change of coordinates $z=F(z')$, $v|dz|^t$ pulls back to $(v\circ
F)|\det F'|^t|dz'|^t$.  Smooth compactly supported densities can be integrated,
and $\CmI(X;\Omegah)$ is the dual of $\CIc(X;\Omegah)$, the pairing of two
half-densities being a density; locally an element of $\CmI(X;\Omegah)$ is
$v|dz|^{1/2}$ with $v$ a distribution in coordinates, transforming as above.
A smooth vector field $V$ acts on half-densities by the Lie derivative: in
coordinates, with $V=\sum v_j\pa_{z_j}$,
\begin{equation}\label{conormal:eq:Lie}
  \mathcal L_V(\phi|dz|^{t})=\big(V\phi+t(\textstyle\sum_j\pa_{z_j}v_j)\phi\big)|dz|^{t}.
\end{equation}

Let $Y\subset X$ be a closed embedded submanifold of codimension $k$.  An
\emph{adapted chart} is a chart $\psi:U\to\psi(U)\subset\R^n=\R^k_x\times\R^d_y$
with $\psi(U\cap Y)=\psi(U)\cap\{x=0\}$; such charts cover $Y$.  The
\emph{normal bundle} $NY$ has fibres $N_yY=T_yX/T_yY$, and the \emph{conormal
bundle} $N^*Y\subset T^*X$, with fibres $N_y^*Y=\{\zeta\in T^*_yX;\
\zeta(T_yY)=0\}$, is its dual.  In an adapted chart,
$N^*Y=\{(0,y;\xi,0)\}$, and $(y,\xi)$ are coordinates on $N^*Y$.  If
$F:\Omega'\to\Omega$ is a change of adapted charts as in
\cref{conormal:sec:invariance}, the pull-back
$F^*(\xi\cdot dx)=(A(y)^t\xi)\cdot dx'$ at $(0,y)$, since $\pa_yF'(0,y)=0$; so
the induced identification of conormal bundles is
\begin{equation}\label{conormal:eq:Fsharp}
  F^\sharp:N^*Y\ni(y,\eta)\longmapsto\big(g(y),A(y)^{-t}\eta\big)\in N^*Y,
\end{equation}
(primed coordinates on the left, unprimed on the right) which is linear on the
fibres.

\begin{definition}\label{conormal:def:manifold}
  The space $\Ic{m}(X,Y;\Omegah)$ consists of those $u\in\CmI(X;\Omegah)$ which
  are smooth on $X\setminus Y$ and such that for every adapted chart
  $\psi:U\to\psi(U)$ and every $\phi\in\CIc(U)$,
  $(\psi^{-1})^*(\phi u)=v|dz|^{1/2}$ with $v\in\Ic{m}_{\mathrm c}(\psi(U),\{x=0\})$.
  The classical subspace $\Ic{m}_{\mathrm{ph}}(X,Y;\Omegah)$ is defined in the
  same way.  For a vector bundle $E$ over $X$, $\Ic{m}(X,Y;E\otimes\Omegah)$ is
  defined by requiring the coefficients in local trivializations to be in
  $\Ic{m}(X,Y;\Omegah)$; similarly with $\Omegah$ omitted.
\end{definition}

By \cref{conormal:thm:invariance} (the factor $|\det F'|^{1/2}$ being smooth)
it suffices to check the condition for a family of adapted charts covering $Y$,
and the definition agrees with \cref{conormal:def:local} for $X=\R^n$ and with
\cref{conormal:def:bundle} for the zero section of a vector bundle; it is
independent of the choice of transition functions for $E$ by the module
property.

\subsection{Normal fibrations}
Alternatively, $\Ic{m}(X,Y)$ can be defined by transferring
\cref{conormal:def:bundle} from the normal bundle to $X$
(\cref{conormal:fig:tubular}).

\begin{figure}[htbp]
\centering
\begin{tikzpicture}[>=Latex,scale=0.9]
  \begin{scope}
    \fill[gray!15] plot[smooth] coordinates {(-2.6,-0.2) (-1.3,0.95) (0,0.75) (1.3,1.25) (2.6,0.75)}
       -- plot[smooth] coordinates {(2.6,-0.75) (1.3,-0.25) (0,-0.75) (-1.3,-0.55) (-2.6,-1.4)} -- cycle;
    \draw[very thick] plot[smooth] coordinates {(-2.6,-0.8) (-1.3,0.2) (0,0) (1.3,0.5) (2.6,0)};
    \foreach \p/\q in {(-1.95,-0.33)/(-1.55,-0.95),(-1.3,0.2)/(-1.3,-0.55),(-0.65,0.17)/(-0.6,-0.6),(0,0)/(0,-0.75),
                       (0.65,0.2)/(0.62,-0.53),(1.3,0.5)/(1.3,-0.25),(1.95,0.33)/(1.97,-0.43)}
      \draw[gray] \p--\q;
    \foreach \p/\q in {(-1.95,-0.33)/(-2.25,0.25),(-1.3,0.2)/(-1.3,0.95),(-0.65,0.17)/(-0.7,0.92),(0,0)/(0,0.75),
                       (0.65,0.2)/(0.68,1.0),(1.3,0.5)/(1.3,1.25),(1.95,0.33)/(1.93,1.08)}
      \draw[gray] \p--\q;
    \node[font=\footnotesize] at (2.75,0.2) {$Y$};
    \node[font=\footnotesize] at (-2.3,1.2) {$D\subset X$};
  \end{scope}
  \draw[->] (3.2,0)--(4.4,0) node[midway,above,font=\footnotesize]{$f$};
  \begin{scope}[xshift=7.4cm]
    \fill[gray!15] (-2.6,-0.75) rectangle (2.6,0.75);
    \draw[very thick] (-2.6,0)--(2.6,0);
    \foreach \x in {-1.95,-1.3,...,1.95} \draw[gray] (\x,-0.75)--(\x,0.75);
    \node[font=\footnotesize] at (3.05,0) {$0_{NY}$};
    \node[font=\footnotesize] at (-2.1,1.05) {$D'\subset NY$};
    \node[font=\footnotesize] at (1.62,-1.05) {$N_yY$};
  \end{scope}
\end{tikzpicture}
\caption[A normal fibration]{A normal fibration
(\cref{conormal:thm:tubular}) identifies a neighbourhood $D$ of $Y$ in $X$
with a neighbourhood $D'$ of the zero section in the normal bundle; the
gray curves on the left (for instance normal geodesics) go to the fibres
$N_yY$. Since its differential along $Y$ induces the identity on $N_yY$, two
choices differ by a map of the type in
\cref{conormal:prop:bundle-invariance}, which preserves conormal
distributions and their principal symbols.}
\label{conormal:fig:tubular}
\end{figure}

\begin{theorem}[Tubular neighbourhood theorem]\label{conormal:thm:tubular}
  If $Y\subset X$ is a closed embedded submanifold, there are an open
  neighbourhood $D$ of $Y$ in $X$, an open neighbourhood $D'$ of the zero section
  of $NY$, and a diffeomorphism $f:D\to D'$ (a \emph{normal fibration}) such
  that $f|_Y$ is the natural identification of $Y$ with $0_{NY}$ and, for each
  $y\in Y$, the map induced by $f_*$ from $N_yY=T_yX/T_yY$ to
  $T_y(NY)/T_y0_{NY}=N_yY$ is the identity.  If $f_1,f_2$ are two normal
  fibrations then $f_2\circ f_1^{-1}$ is a diffeomorphism between neighbourhoods
  of $0_{NY}$ satisfying the hypotheses of
  \cref{conormal:prop:bundle-invariance}.
\end{theorem}
\begin{proof}[Sketch of proof]
  Choose a Riemannian metric on $X$; it identifies $N_yY$ with the orthogonal
  complement of $T_yY$ in $T_yX$.  The exponential map $NY\ni(y,v)\mapsto\exp_y(v)$
  is smooth, is the identity on the zero section and has differential the
  identity at points of the zero section (with respect to the splitting
  $T_y(NY)=T_yY\oplus N_yY$), so it is a local diffeomorphism near $0_{NY}$; a
  standard argument (using that it is injective on $0_{NY}$ and $Y$ is closed)
  shows it is a diffeomorphism of a neighbourhood $D'$ of $0_{NY}$ onto a
  neighbourhood $D$ of $Y$, and $f$ is its inverse.  The last statement is
  immediate from the definition.
\end{proof}

\begin{proposition}\label{conormal:prop:fibration-def}
  Let $f:D\to D'$ be a normal fibration and $Y$ compact.  Then $u\in\CmI(X)$
  lies in $\Ic{m}(X,Y)$ if and only if $u$ is smooth outside $Y$ and, for some
  (hence any) $\phi\in\CIc(D)$ equal to one near $Y$, $(f^{-1})^*(\phi u)\in
  \Ic{m}(NY,0_{NY})$.
\end{proposition}
\begin{proof}
  Local trivializations of $NY$ composed with $f$ are adapted charts for $Y$, and
  local trivializations of $NY$ are adapted charts for $0_{NY}$, so both
  conditions amount to the condition of \cref{conormal:def:manifold} in a
  particular covering family of adapted charts.  Independence of $f$ also
  follows from \cref{conormal:prop:bundle-invariance} and
  \cref{conormal:thm:tubular}.
\end{proof}

\subsection{The principal symbol}
\begin{discussion}
  Why should the principal symbol be a half-density? Under a change of
  adapted coordinates the full symbol picks up the factor
  $\abs{\det A(y)}^{-1}$ from the Fourier transform in the normal variables
  (\cref{conormal:thm:invariance}). If $u$ itself is a half-density, the
  Jacobian factor $\abs{\det F'}^{1/2}$ of $u$ combines with it to give
  exactly the transformation law of a half-density on $N^*Y$ in the
  coordinates $(y,\xi)$. With this convention the principal symbol is
  intrinsic, with no choices of densities on $X$ or $Y$; this is also the
  form needed for Lagrangian distributions in \cref{ch:hj}.
\end{discussion}

For $r\in\R$ let $S^r(N^*Y;\Omegah)$ be the space of smooth half-densities on
$N^*Y$ which in each adapted chart are of the form $\beta(y,\xi)|dy\,d\xi|^{1/2}$
with $\phi\beta\in S^{r-\frac k2}(\R^d;\R^k)$ for all $\phi\in\CIc$ of the chart.
This is invariant under the fibre-linear maps \eqref{conormal:eq:Fsharp}, so is
well defined; the shift by $k/2$ reflects the fact that $|d\xi|^{1/2}$ is
homogeneous of degree $k/2$ under fibre dilations, so that the order $r$ is the
total homogeneity.  Write
$S^{r-[1]}(N^*Y;\Omegah)=S^r(N^*Y;\Omegah)/S^{r-1}(N^*Y;\Omegah)$.  In the
classical case the principal parts are smooth half-densities on
$N^*Y\setminus0$ homogeneous of degree $r$; we denote this space
$\CI_{(r)}(N^*Y\setminus0;\Omegah)$.

\begin{theorem}\label{conormal:thm:symbol}
  There is a unique map
  \begin{equation*}
    \sigma_m:\Ic{m}(X,Y;\Omegah)\longrightarrow S^{m+\frac n4-[1]}(N^*Y;\Omegah)
  \end{equation*}
  such that if $u=(\osc{a}+f)|dx\,dy|^{1/2}$ in an adapted chart then
  $\sigma_m(u)=a(y,\xi)|dy\,d\xi|^{1/2}$ there.  It gives a short exact sequence
  \begin{equation}\label{conormal:eq:exact}
    0\longrightarrow\Ic{m-1}(X,Y;\Omegah)\longrightarrow\Ic{m}(X,Y;\Omegah)
    \xrightarrow{\ \sigma_m\ }S^{m+\frac n4-[1]}(N^*Y;\Omegah)\longrightarrow0,
  \end{equation}
  and, in the classical case, a short exact sequence
  \begin{equation*}
    0\to\Ic{m-1}_{\mathrm{ph}}(X,Y;\Omegah)\to\Ic{m}_{\mathrm{ph}}(X,Y;\Omegah)
    \xrightarrow{\sigma_m}\CI_{(m+\frac n4)}(N^*Y\setminus0;\Omegah)\to0.
  \end{equation*}
  Moreover $\sigma_m(\psi u)=(\psi|_Y)\sigma_m(u)$ for $\psi\in\CI(X)$, where
  $\psi|_Y$ is lifted to $N^*Y$.
\end{theorem}
\begin{proof}
  We must check that the local definition is consistent.  Let $F$ be a change of
  adapted coordinates.  If $u=v|dz|^{1/2}$ then $F^*u=(F^*v)|\det
  F'|^{1/2}|dz'|^{1/2}$.  At $x=0$, $F'$ is block triangular with diagonal
  blocks $A(y)$ and $g'(y)$, so $|\det F'(0,y)|^{1/2}=|\det A|^{1/2}|\det
  g'|^{1/2}$; by \cref{conormal:thm:invariance} and the multiplication rule in
  \cref{conormal:prop:local-basic}(v), the full symbol of $F^*u$ is congruent
  modulo $S^{M-1}$ to
  \begin{equation*}
    |\det A(y)|^{-\frac12}|\det g'(y)|^{\frac12}a\big(g(y),A(y)^{-t}\eta\big).
  \end{equation*}
  On the other hand the Jacobian determinant of $F^\sharp$ in
  \eqref{conormal:eq:Fsharp} is $\det g'(y)\det A(y)^{-1}$, so
  $(F^\sharp)^*(a|dy\,d\xi|^{1/2})$ is exactly this function times
  $|dy\,d\eta|^{1/2}$.  Hence $\sigma_m$ is well defined, with
  $\sigma_m(F^*u)=(F^\sharp)^*\sigma_m(u)$; the classical case is the same.
  The multiplicative property holds because $\psi(x,y)-\psi(0,y)=\sum_jx_j\psi_j$
  and $x_j$ lowers the order.  Exactness in the middle is local and is
  \cref{conormal:prop:local-basic}(iii); injectivity is trivial.  For
  surjectivity, given $b\in S^{m+\frac n4}(N^*Y;\Omegah)$, cover $Y$ by adapted
  charts $U_j$ (a locally finite cover) with a partition of unity $\phi_j$ on $Y$,
  write $\phi_jb=\beta_j|dy\,d\xi|^{1/2}$ and let
  $u=\sum_j\chi_j\osc{\beta_j}|dz|^{1/2}$, with $\chi_j\in\CIc(U_j)$ equal to one
  near $\supp\phi_j\times\{0\}$; then $\sigma_m(u)=\sum_j\phi_jb=b$.
\end{proof}

\subsection{Properties}
We collect the main properties of the spaces $\Ic{m}(X,Y;\Omegah)$; all hold
equally for the classical spaces.

\begin{theorem}\label{conormal:thm:properties}
  Let $Y\subset X$ be a closed embedded submanifold of codimension $k$, $\dim
  X=n$.
  \begin{enumerate}[label=(\roman*)]
  \item Elements of $\Ic{m}(X,Y;\Omegah)$ are smooth outside $Y$,
    $\Ic{m-1}\subset\Ic{m}$ and $\bigcap_m\Ic{m}(X,Y;\Omegah)=\CI(X;\Omegah)$.
  \item \emph{(Asymptotic completeness)} If $u_j\in\Ic{m-j}(X,Y;\Omegah)$, there
    exists $u\in\Ic{m}(X,Y;\Omegah)$ with $u-\sum_{j<N}u_j\in\Ic{m-N}$ for all $N$.
  \item \emph{(Invariance)} If $F:X'\to X$ is a diffeomorphism and
    $Y'=F^{-1}(Y)$, then $F^*:\Ic{m}(X,Y;\Omegah)\to\Ic{m}(X',Y';\Omegah)$ and
    $\sigma_m(F^*u)=(F^\sharp)^*\sigma_m(u)$, where $F^\sharp:N^*Y'\to N^*Y$ is the
    inverse of the restriction of $F^*:T^*X\to T^*X'$.  The same holds for
    diffeomorphisms between open subsets, for distributions supported in the
    range.
  \item \emph{(Differential operators)} If $P\in\Diff{q}(X;\Omegah)$ then
    $P:\Ic{m}(X,Y;\Omegah)\to\Ic{m+q}(X,Y;\Omegah)$ and
    $\sigma_{m+q}(Pu)=\sigma_q(P)|_{N^*Y}\,\sigma_m(u)$, where $\sigma_q(P)$ is the
    principal symbol, characterized by $P(e^{itf}v)=e^{itf}(t^q\sigma_q(P)(df)v
    +O(t^{q-1}))$ as $t\to\infty$ for real $f\in\CI(X)$ and
    $v\in\CIc(X;\Omegah)$.
  \item \emph{(Wavefront set)} $\WF(u)\subset N^*Y\setminus0$ for
    $u\in\Ic{m}(X,Y;\Omegah)$.
  \end{enumerate}
\end{theorem}
\begin{proof}
  All statements are local and invariant, so follow from the local results:
  (i) from \cref{conormal:prop:local-basic}(ii); (ii) from
  \cref{conormal:prop:local-basic}(iv) and a partition of unity; (iii) from
  \cref{conormal:thm:invariance,conormal:thm:symbol}; (iv) from
  \cref{conormal:prop:local-basic}(v), since in coordinates
  $\sigma_q(P)(z,\zeta)=p_q(z,\zeta)$ (the half-density factor and the lower order
  terms only contribute terms of order $q-1$); (v) from
  \cref{conormal:prop:WF-local}.
\end{proof}

\begin{remark}\label{conormal:rem:scalar}
  If one works with functions rather than half-densities, the computation in the
  proof of \cref{conormal:thm:symbol} shows that the full symbol transforms by
  \eqref{conormal:eq:symbol-transform}, i.e.\ $a|d\xi|$ transforms as a fibre
  density on $N^*Y$: the principal symbol of $u\in\Ic{m}(X,Y)$ is then a
  section of $\Omega_{\mathrm{fib}}(N^*Y)$ with coefficient of order $M$, as in
  \eqref{conormal:eq:bundle-exact}.  Half-densities make the symbol a
  half-density on the symplectic manifold $N^*Y$, which is the form used in
  \cref{ch:hj}.
\end{remark}

\section{Iterated regularity}\label{conormal:sec:iterated}

We now show that conormal distributions are characterized by a regularity
property which is manifestly invariant and does not involve the Fourier
transform.  The idea is that differentiating along $Y$, or applying
$x_i\pa_{x_j}$, which is homogeneous of degree $0$ under dilations normal to $Y$,
does not make a conormal singularity worse
(\cref{conormal:fig:local}(a)), whereas a general derivative raises the order
by one.  Let $\Vt(X,Y)$ denote the space of smooth vector fields on $X$
tangent to $Y$; it is a Lie algebra and a $\CI(X)$-module.

\begin{lemma}\label{conormal:lem:generators}
  In the local model, $\Vt(\Omega,Y)$ is spanned over $\CI(\Omega)$ by the
  vector fields $x_i\pa_{x_j}$, $1\le i,j\le k$, and $\pa_{y_l}$, $1\le l\le d$.
\end{lemma}
\begin{proof}
  A vector field $V=\sum_ja_j\pa_{x_j}+\sum_lb_l\pa_{y_l}$ is tangent to
  $\{x=0\}$ if and only if $a_j(0,y)=0$ for all $j$, that is, $a_j=\sum_i
  x_ia_{ij}$ with $a_{ij}\in\CI$ (Taylor's formula with integral remainder in $x$,
  valid in a neighbourhood of each point of $Y$; away from $Y$ use
  $\pa_{x_j}=\sum_i(x_i/|x|^2)x_i\pa_{x_j}$ and a partition of unity).
\end{proof}

\begin{definition}\label{conormal:def:IH}
  For $s\in\R$ let
  \begin{equation*}
    \IH{s}(X,Y;\Omegah)=\{u\in H^s_{\mathrm{loc}}(X;\Omegah);\
    \mathcal L_{V_1}\cdots\mathcal L_{V_j}u\in H^s_{\mathrm{loc}}(X;\Omegah)\
    \forall\ j,\ V_i\in\Vt(X,Y)\},
  \end{equation*}
  and $\Ic{*}(X,Y;\Omegah)=\bigcup_s\IH{s}(X,Y;\Omegah)$.
\end{definition}
Here $H^s_{\mathrm{loc}}$ is defined via coordinates and is invariant under
diffeomorphisms (\cref{ch:microlocal}); the zeroth order terms in
\eqref{conormal:eq:Lie} are irrelevant since $H^s_{\mathrm{loc}}$ is a
$\CI$-module, so in coordinates we may replace $\mathcal L_V$ by $V$ acting on
coefficients.  Thus $\IH{s}$ is manifestly invariant under diffeomorphisms
preserving $Y$, and $\bigcap_s\IH{s}(X,Y;\Omegah)=\CI(X;\Omegah)$ by Sobolev
embedding.

\begin{lemma}\label{conormal:lem:IH-local}
  Let $u\in H^s(\R^n)$ have compact support.  Then $u\in\IH{s}(\R^n,Y)$ if and
  only if
  \begin{equation}\label{conormal:eq:IH-local}
    x^\alpha D_x^\beta D_y^\gamma u\in H^s(\R^n)\quad\text{for all }
    \alpha,\beta,\gamma\text{ with }|\beta|\le|\alpha|,
  \end{equation}
  or equivalently if and only if $D_x^\beta D_y^\gamma x^\alpha u\in H^s$ for
  all $|\beta|\le|\alpha|$.
\end{lemma}
\begin{proof}
  If $V_i\in\Vt$, then by \cref{conormal:lem:generators}, moving coefficients to
  the left using $[W,\psi]=W\psi$, $V_1\cdots V_ju$ is a combination with smooth
  coefficients of products of at most $j$ of the generators $x_i\pa_{x_j}$,
  $\pa_{y_l}$ applied to $u$; and such products are linear combinations with
  constant coefficients of $x^\alpha D_x^\beta D_y^\gamma$ with $|\alpha|=|\beta|$.
  Conversely, by induction on $|\alpha|$, $x^\alpha D_x^\beta$ with
  $|\alpha|=|\beta|$ is a combination of products of the $x_iD_{x_j}$, since
  $x_iD_{x_j}x^{\alpha'}D_x^{\beta'}=x_ix^{\alpha'}D_{x_j}D_x^{\beta'}+
  x_i(D_{x_j}x^{\alpha'})D_x^{\beta'}$ and the last term is shorter.  The cases
  $|\beta|<|\alpha|$ follow by multiplying by $\chi x^{\alpha''}$, with
  $\chi\in\CIc$ equal to one on $\supp u$; and $D^\beta x^\alpha$ is a
  combination of $x^{\alpha-\delta}D^{\beta-\delta}$, which preserves the
  condition $|\beta|\le|\alpha|$.
\end{proof}

The key analytic step converts weighted $L^2$ bounds on all derivatives into
symbol estimates.

\begin{lemma}\label{conormal:lem:L2-sup}
  Let $t\in\R$, $j\in\N_0$, $L>j+\frac k2$, and let $b\in\CI(\R^k)$ satisfy
  $\ang\xi^{t+|\alpha|}D^\alpha b\in L^2(\R^k)$ for $|\alpha|\le L$.  Then
  \begin{equation*}
    |D^\beta b(\xi)|\le C\ang\xi^{-t-\frac k2-|\beta|},\quad |\beta|\le j,
    \qquad C=C_{k,j,t,L}\sum_{|\alpha|\le L}\norm{\ang\xi^{t+|\alpha|}D^\alpha b}_{L^2}.
  \end{equation*}
  Conversely, if $b\in S^\mu(\R^k)$ then $\ang\xi^{t+|\alpha|}D^\alpha b\in L^2$
  for all $\alpha$ provided $t<-\mu-\frac k2$.
\end{lemma}
\begin{proof}
  The converse is immediate since $\ang\xi^{t+\mu}\in L^2$ if $2(t+\mu)<-k$.
  For the direct statement, let $A=\{\frac12<|\zeta|<2\}$ and $\chi\in\CIc(A)$
  equal to one on $|\zeta|=1$.  For $R\ge1$ set $b_R(\zeta)=b(R\zeta)$.  On
  $RA$ we have $R/2\le\ang\xi\le3R$, so
  \begin{equation*}
    \norm{D^\alpha b_R}^2_{L^2(A)}=R^{2|\alpha|-k}\int_{RA}|D^\alpha b|^2d\xi
    \le CR^{2|\alpha|-k}R^{-2(t+|\alpha|)}\norm{\ang\xi^{t+|\alpha|}D^\alpha b}^2_{L^2}.
  \end{equation*}
  Hence $\norm{\chi b_R}_{H^L}\le CR^{-t-\frac k2}$, and by Sobolev's embedding
  $|D^\beta b_R(\zeta)|\le CR^{-t-\frac k2}$ for $|\zeta|=1$, $|\beta|\le j$.
  Given $|\xi|\ge1$ take $R=|\xi|$, $\zeta=\xi/|\xi|$: then $D^\beta b(\xi)=
  R^{-|\beta|}(D^\beta b_R)(\zeta)$, which gives the estimate.  For $|\xi|\le1$
  apply Sobolev's embedding to $\chi_0b$ with $\chi_0\in\CIc$ equal to one on the
  unit ball.
\end{proof}

\begin{theorem}\label{conormal:thm:IH}
  In the local model, for compactly supported distributions,
  \begin{equation}\label{conormal:eq:IH-incl}
    \Ic{m}_{\mathrm c}(\R^n,Y)\subset\IH{s}(\R^n,Y)\ \text{ if } s<-m-\frac n4,
    \qquad
    \IH{s}(\R^n,Y)\cap\CmIc(\R^n)\subset\Ic{-s-\frac n4}_{\mathrm c}(\R^n,Y).
  \end{equation}
  Consequently, for any closed embedded submanifold $Y\subset X$,
  \begin{equation*}
    \IH{s}(X,Y;\Omegah)\subset\Ic{-s-\frac n4}(X,Y;\Omegah)\subset
    \IH{s'}(X,Y;\Omegah)\quad\forall\ s'<s,
  \end{equation*}
  so $\Ic{*}(X,Y;\Omegah)=\bigcup_m\Ic{m}(X,Y;\Omegah)$.
\end{theorem}
\begin{proof}
  First inclusion.  Let $u=\osc{c}+f$ with $c\in S^M$ compactly supported in
  $(x,y)$ and $f\in\CIc$.  By \cref{conormal:lem:osc-ops}(i), each
  $x^\alpha D_x^\beta D_y^\gamma$ with $|\alpha|=|\beta|$ maps $\osc{c}$ to some
  $\osc{c'}$ with $c'$ of order $M$ and compact support, so by
  \cref{conormal:lem:IH-local} it suffices to show $\osc{c}\in H^s$.  Let
  $\hat c(\zeta',\eta;\xi)$ be the Fourier transform of $c(\cdot,\cdot,\xi)$ in
  $(x,y)$; since $c$ has compact support in $(x,y)$,
  $|\hat c(\zeta',\eta;\xi)|\le C_N\ang{\zeta'}^{-N}\ang{\eta}^{-N}\ang\xi^{M}$.
  Then (first for $M<-k$, by Fubini, then in general by
  \cref{conormal:lem:density})
  \begin{align*}
    \widehat{\osc{c}}(\zeta,\eta)&=\int\hat c(\zeta-\xi,\eta;\xi)\,\dbar\xi,\\
    |\widehat{\osc{c}}(\zeta,\eta)|&\le C_N\ang\eta^{-N}\int
    \ang{\zeta-\xi}^{-N}\ang\xi^M\,d\xi\le C'_N\ang\eta^{-N}\ang\zeta^M
  \end{align*}
  by Peetre's inequality, for $N>k+|M|$.  Since
  $\ang{(\zeta,\eta)}^{s}\le\ang\zeta^s\ang\eta^{|s|}$,
  $\ang{(\zeta,\eta)}^s\widehat{\osc{c}}\in L^2$ provided $2(s+M)<-k$, i.e.\
  $s<-M-\frac k2=-m-\frac n4$.

  Second inclusion.  Let $u\in\IH{s}$ have compact support, so $\hat u$ is smooth.
  By \cref{conormal:lem:IH-local} and $\cF(x^\alpha u)=(-D_\zeta)^\alpha\hat u$,
  $\zeta^\beta\eta^\gamma D_\zeta^\alpha\hat u\in\ang{(\zeta,\eta)}^{-s}L^2$ for
  $|\beta|\le|\alpha|$ and all $\gamma$.  Since $\ang\zeta^{|\alpha|}\le
  C\sum_{|\beta|\le|\alpha|}|\zeta^\beta|$, similarly for $\ang\eta^N$, and
  $\ang{(\zeta,\eta)}^{s}\ge c\ang\zeta^{s}\ang\eta^{-|s|}$, we get
  \begin{equation}\label{conormal:eq:weighted}
    \ang\zeta^{s+|\alpha|}\ang\eta^{N}D_\zeta^\alpha\hat u(\zeta,\eta)\in
    L^2(\R^n)\quad\forall\ \alpha,N.
  \end{equation}
  Let $a(y,\zeta)=\int e^{iy\cdot\eta}\hat u(\zeta,\eta)\,\dbar\eta$ be the
  partial Fourier transform of $u$ in $x$; it is smooth, has compact support in
  $y$, and $u=\osc{a}$.  By Minkowski's and the Cauchy--Schwarz inequalities, for
  $N>|\gamma|+\frac d2$,
  \begin{equation*}
    \norm{\ang\zeta^{s+|\alpha|}D_\zeta^\alpha D_y^\gamma a(y,\cdot)}_{L^2}\le
    \int|\eta|^{|\gamma|}\norm{\ang\zeta^{s+|\alpha|}D_\zeta^\alpha
    \hat u(\cdot,\eta)}_{L^2}\dbar\eta
    \le C\norm{\ang\eta^N\ang\zeta^{s+|\alpha|}D_\zeta^\alpha\hat u}_{L^2(\R^n)},
  \end{equation*}
  uniformly in $y$.  Applying \cref{conormal:lem:L2-sup} with $t=s$ to
  $D_y^\gamma a(y,\cdot)$ gives $|D_\zeta^\beta D_y^\gamma a|\le
  C\ang\zeta^{-s-\frac k2-|\beta|}$, so $a\in S^{-s-\frac k2}_{\mathrm c}$, which
  by \eqref{conormal:eq:order} means $u\in\Ic{-s-\frac n4}_{\mathrm c}$.

  The global statement follows by localization in adapted charts, using the
  invariance of $\IH{s}$ and of $\Ic{m}$.  Note that the middle space in the
  global chain is $\Ic{-s-\frac n4}$, and $\Ic{-s-\frac n4}\subset\IH{s'}$ for
  $s'<s$ by the first inclusion.
\end{proof}

\begin{question}
  Why do the two inclusions in \cref{conormal:thm:IH} not match exactly?
\end{question}
\begin{answer}
  Sobolev spaces measure $L^2$ norms, symbol classes measure supremum norms,
  and $L^2$ cannot see a borderline power exactly. Already for $\delta_0$, of
  order $\frac n4$, we have $\delta_0\in H^s$ if and only if $s<-\frac n2$,
  since $\ang\zeta^{s}\in L^2(\R^n)$ if and only if $s<-\frac n2$. So
  $\delta_0\in\IH{s}$ for $s<-\frac n2=-\frac n4-\frac n4$ but not for
  $s=-\frac n2$. Conversely the passage from weighted $L^2$ bounds to
  pointwise bounds in \cref{conormal:lem:L2-sup} costs exactly $\frac k2$ in
  the order. The net effect is an arbitrarily small loss.
\end{answer}

\begin{remark}
  \Cref{conormal:thm:IH} gives a second proof that $\Ic{*}(X,Y)$ is invariant
  under diffeomorphisms, and that the order is invariant up to an arbitrarily
  small loss; the sharp invariance of the order and of the principal symbol
  requires \cref{conormal:thm:invariance}.  For $Y=\{0\}$ the theorem is the
  statement that $D_\zeta^\alpha\hat u\in\ang\zeta^{-s-|\alpha|}L^2$ for all
  $\alpha$ implies $\hat u\in S^{-s-\frac n2}$, and conversely
  $S^{\mu}\subset\ang\zeta^{-s-|\alpha|}L^2$-type spaces for $s<-\mu-\frac n2$.
\end{remark}

\section{Examples}\label{conormal:sec:examples}

\begin{example}[Delta sections]\label{conormal:ex:delta}
  In the local model, $\delta(x)=\delta(x_1)\cdots\delta(x_k)=\osc{1}$, the full
  symbol being $1\in S^0$; so $\delta(x)\in\Ic{m}_{\mathrm{ph}}(\R^n,Y)$ with
  $M=0$, i.e.\ $m=\frac k2-\frac n4=\frac14(\operatorname{codim}Y-\dim Y)$.
  Invariantly: the restriction to $Y$ of the density bundle of $X$ splits as
  $\Omega_YX=\Omega Y\otimes\Omega(NY)$, where $\Omega(NY)$ is the bundle of
  densities on the fibres of the normal bundle.  If $0<\mu\in\CI(Y;\Omega(NY))$
  is a positive normal density, then
  \begin{equation*}
    u_\mu:\CIc(X;\Omega X)\ni\nu\longmapsto\int_Y\nu|_Y/\mu
  \end{equation*}
  is a well-defined distribution, and in adapted coordinates in which
  $\mu=\mu(y)|dx|$ it is $\mu(y)^{-1}\delta(x)$.  So $u_\mu\in
  \Ic{\frac14(k-d)}_{\mathrm{ph}}(X,Y)$ with full symbol $\mu(y)^{-1}$; with the
  normalization of \cref{conormal:rem:scalar}, its principal symbol is
  $\mu(y)^{-1}|d\xi|$, the fibre density on $N^*Y$ dual to $\mu$.
\end{example}

\begin{example}[Poincar\'e dual forms]\label{conormal:ex:poincare}
  Conormal sections of vector bundles are defined in
  \cref{conormal:def:manifold}.  Let $X$ be a compact oriented manifold of
  dimension $n$ and $Y\subset X$ a closed embedded oriented submanifold of
  dimension $d$.  Near any point of $Y$ choose oriented coordinates in which
  $Y=\{z_{d+1}=\dots=z_n=0\}$ and $(z_1,\dots,z_d)$ are oriented coordinates on $Y$,
  and set
  \begin{equation*}
    p_Y=\delta(z_{d+1})\cdots\delta(z_n)\,dz_{d+1}\wedge\dots\wedge dz_n\in
    \Ic{\frac14(n-2d)}_{\mathrm{ph}}(X,Y;\Lambda^{n-d}).
  \end{equation*}
  This is independent of the choices, because it is characterized by
  \begin{equation*}
    \int_X\phi\wedge p_Y=\int_Y\phi\qquad\forall\ \phi\in\CI(X;\Lambda^d),
  \end{equation*}
  as one checks in the given coordinates ($\phi\wedge p_Y$ only involves the
  coefficient of $dz_1\wedge\dots\wedge dz_d$ in $\phi$).  It follows that
  $dp_Y=0$ in the distributional sense: for $\psi\in\CI(X;\Lambda^{d-1})$,
  $\int_X\psi\wedge dp_Y=\pm\int_Xd\psi\wedge p_Y=\pm\int_Yd\psi=0$ by Stokes'
  theorem.  The distributional de Rham cohomology, defined by the complex
  $d:\CmI(X;\Lambda^j)\to\CmI(X;\Lambda^{j+1})$, is canonically isomorphic to the
  smooth de Rham cohomology (this follows, for instance, from the Hodge theorem),
  and the class of $p_Y$ in $H^{n-d}_{\mathrm{dR}}(X)$ is, by the displayed
  identity, the Poincar\'e dual of $Y$.
\end{example}

\section{Exercises}\label{conormal:sec:exercises}

\begin{exercise}\label{conormal:ex:Vkl}
  Show by induction that $a\in S^0(\R^n)$ if and only if $a\in\CI(\R^n)$ and
  $V_1\cdots V_Na\in L^\infty(\R^n)$ for all $N$ and all choices of
  $V_j\in\{\xi_k\pa_{\xi_l}\}$ (including $N=0$).  (Hint: in $|\xi|\ge1$,
  $\pa_{\xi_l}=\sum_k\xi_k|\xi|^{-2}\,\xi_k\pa_{\xi_l}$.)
\end{exercise}

\begin{exercise}\label{conormal:ex:Ab}
  For a compact manifold with boundary $X$ let $\mathcal V_{\mathrm b}(X)$ be the
  Lie algebra of smooth vector fields tangent to $\pa X$ and
  \begin{equation*}
    \mathcal A(X)=\{a\in\CI(X\setminus\pa X)\cap L^\infty;\ V_1\cdots V_Na\in
    L^\infty\ \forall\ N,\ V_i\in\mathcal V_{\mathrm b}(X)\}.
  \end{equation*}
  Using that the linear vector fields $\xi_k\pa_{\xi_l}$ span
  $\mathcal V_{\mathrm b}(\RC{\R^n})$ near the boundary, show that
  $S^m(\R^n)=\rho^{-m}\mathcal A(\RC{\R^n})$.
\end{exercise}

\begin{exercise}\label{conormal:ex:blowup}
  Let $W$ be a real vector space of dimension $n$.  Define the sphere of $W$ as
  $\Sphere W=(W\setminus0)/\R^+$.  Show that $[W,\{0\}]=(W\setminus0)\sqcup\Sphere W$
  has a unique structure of manifold with boundary, diffeomorphic to
  $[0,\infty)\times\Sphere^{n-1}$, such that the blow-down map $\beta:[W,\{0\}]\to
  W$ (the identity on $W\setminus0$, mapping $\Sphere W$ to $0$) is smooth and the
  action of $\GL(W)$ lifts to a smooth action (use \cref{conormal:lem:blowup}).
\end{exercise}

\begin{exercise}
  Let $U$ be a finite-dimensional complex vector space and set
  $\Ic{m}(W,\{0\};U)=\Ic{m}(W,\{0\})\otimes U$.  Show that the Fourier transform
  gives an isomorphism $\Ic{m}(W,\{0\};U)\to S^{m-\frac n4}(W';U\otimes\Omega W')$
  (with $n=\dim W$), that $\Omega U'\otimes\Omega U$ is canonically trivial for a
  real vector space $U$, and that the integral of
  \cref{conormal:lem:point-basic}(iv) defines a linearly invariant map
  $\Ic{m}(W,\{0\};\Omega W)\to\C$.
\end{exercise}

\begin{exercise}
  Let $u\in\CmI(\R^n)$ be smooth outside $0$ and homogeneous modulo $\CI$ of
  degree $h$: $u(tz)=t^hu(z)+F(t,z)$ for $t>0$, $z\ne0$, with
  $F\in\CI((0,\infty)\times\R^n)$.  Show that $\chi u$ is a classical conormal
  distribution at $0$ for $\chi\in\CIc(\R^n)$, and determine its order.
\end{exercise}

\begin{exercise}
  Show that every $u\in\Ic{m}(X,Y)$ with $\supp u\subset Y$ is of the form $Pu_\mu$
  for some $P\in\Diff{q}(X)$, with $u_\mu$ as in \cref{conormal:ex:delta}, at
  least locally.  Which orders $m$ can occur?
\end{exercise}

\begin{exercise}
  Let $u\in\Ic{m}(X,Y;\Omegah)$ be classical, with principal symbol nonvanishing
  on $N^*Y\setminus0$.  Show that $\WF(u)=N^*Y\setminus0$, and that if
  $A\in\Psido{0}(\R^n)$ satisfies $Au\in\CI$ then the principal symbol of $A$
  vanishes on $N^*Y$ (locally).
\end{exercise}

\section*{Notes}
This chapter merges the treatment of conormal distributions in the 1998 lecture
notes (where they are introduced through iterated regularity) with that of the
2005 lectures (where they are defined by Fourier transform of classical symbols,
first at a point, then at the zero section of a vector bundle, then at a
submanifold via a normal fibration).  The opening argument of the earlier notes,
relating iterated regularity to symbol estimates, has been completed
(\cref{conormal:lem:L2-sup,conormal:thm:IH}), and the invariance under
diffeomorphisms is proved by Moser's homotopy method in the sharp form
$H^*u-u\in\Ic{m-1}$, which yields the invariance of the principal symbol without
the loss $\epsilon$ inherent in the Sobolev approach.  We use H\"ormander's
normalization of the order, for which $\delta_0\in\Ic{n/4}(\R^n,\{0\})$ and the
kernels of pseudodifferential operators of order $m$ have order $m$; some
formulas for the order in the 2005 lectures contain sign slips which have been
corrected.  The standard reference is \cite[Section~18.2]{Hormander-III}; see also
\cite{Hormander-FIO} for the more general Lagrangian distributions, of which
conormal distributions are the special case associated to the Lagrangian
submanifold $N^*Y$.
